\documentclass[11pt]{article}

% ── Packages ──────────────────────────────────────────────────────────────────
\usepackage[T1]{fontenc}
\usepackage[utf8]{inputenc}
\usepackage{lmodern}
\usepackage{microtype}
\usepackage{amsmath, amssymb}
\usepackage{booktabs}
\usepackage{array}
\usepackage{xcolor}
\usepackage{colortbl}
\usepackage{geometry}
\usepackage{titlesec}
\usepackage{enumitem}
\usepackage{parskip}
\usepackage{fancyhdr}
\usepackage{graphicx}
\usepackage{hyperref}

% ── Page layout ───────────────────────────────────────────────────────────────
\geometry{
  letterpaper,
  top    = 1in,
  bottom = 1in,
  left   = 1.1in,
  right  = 1.1in
}

% ── Colors ────────────────────────────────────────────────────────────────────
\definecolor{headerblue}{RGB}{31,  73, 125}
\definecolor{rowblue}   {RGB}{220, 230, 241}
\definecolor{ruleblue}  {RGB}{31,  73, 125}
\definecolor{notesgray} {RGB}{100, 100, 100}
\definecolor{warnred}   {RGB}{180,  30,  30}
\definecolor{greenok}   {RGB}{0,   120,   0}
\definecolor{darkgreen} {RGB}{0,   120,   0}

% ── Section heading styles ────────────────────────────────────────────────────
\titleformat{\section}
  {\large\bfseries\color{headerblue}}
  {\thesection.}{0.6em}{}[\vspace{-4pt}\textcolor{ruleblue}{\rule{\linewidth}{0.8pt}}]

\titleformat{\subsection}
  {\normalsize\bfseries\color{headerblue}}
  {\thesubsection}{0.5em}{}

\titleformat{\subsubsection}
  {\normalsize\itshape\color{headerblue}}
  {\thesubsubsection}{0.5em}{}

\titlespacing*{\section}      {0pt}{18pt}{6pt}
\titlespacing*{\subsection}   {0pt}{12pt}{4pt}
\titlespacing*{\subsubsection}{0pt}{8pt}{3pt}

% ── Header / footer ───────────────────────────────────────────────────────────
\pagestyle{fancy}
\fancyhf{}
\renewcommand{\headrulewidth}{0.4pt}
\fancyhead[L]{\small\textcolor{notesgray}{GNN2 Project --- Speed Benchmark, Impossibility Analysis \& Literature Review}}
\fancyhead[R]{\small\textcolor{notesgray}{March 2026}}
\fancyfoot[C]{\small\thepage}

% ── List spacing ──────────────────────────────────────────────────────────────
\setlist[itemize]{itemsep=3pt, topsep=4pt, leftmargin=1.4em}
\setlist[description]{itemsep=4pt, topsep=4pt, leftmargin=1.6em,
                       labelwidth=0pt, leftmargin=!}

% ── Hyperref ──────────────────────────────────────────────────────────────────
\hypersetup{
  colorlinks = true,
  linkcolor  = headerblue,
  urlcolor   = headerblue,
  pdftitle   = {GNN2 Speed Benchmark and Literature Review},
  pdfauthor  = {GNN2 Project},
}

% ─────────────────────────────────────────────────────────────────────────────
\begin{document}

% ── Title block ───────────────────────────────────────────────────────────────
\begin{center}
  {\LARGE\bfseries\color{headerblue}
    GNN2 Speed Benchmark: OpenDSS vs.\ GNN Models\\[4pt]
    with Impossibility Analysis, Literature Context, and PowerFlowNet Comparison}\\[6pt]
  {\large GNN2 Project \quad|\quad March 2026}
\end{center}

\vspace{4pt}
\noindent\textcolor{ruleblue}{\rule{\linewidth}{1.2pt}}
\vspace{10pt}

\tableofcontents
\vspace{14pt}
\noindent\textcolor{ruleblue}{\rule{\linewidth}{0.6pt}}
\vspace{10pt}

% =============================================================================
\section{Executive Summary}
% =============================================================================

This report combines two investigations into the runtime performance of Graph Neural
Network (GNN) surrogates for distribution power flow.  The first benchmarks the
24-hour (288-timestep) inference speed of two GNN models against OpenDSS on an
95-node distribution feeder and situates those results within the published literature,
including a detailed comparison against PowerFlowMultiNet.  The second presents a
convergent four-layer argument---drawn from benchmark data, framework profiling,
FLOP-count analysis, and direct empirical evidence---establishing that no GNN can
beat an optimised classical solver at CPU batch size~1 in single-scenario sequential
mode.  Together, they give a complete picture of where GNN surrogates are and are not
competitive with traditional power-flow solvers.

\medskip
\noindent\textbf{Core benchmark finding:}
On CPU, OpenDSS is roughly $11.7\times$ faster per timestep than GNN inference.
On GPU with batching, the GNN pipeline comes within $1.6\times$ of OpenDSS.
The GNN's advantage lies in high-throughput, batched, parallel scenario
evaluation---not in single-scenario latency.

\medskip
\noindent\textbf{Impossibility finding:}
No published paper (2019--2026) demonstrates a GNN or any ML surrogate
outperforming an optimised classical solver under the simultaneous conditions
of CPU-only inference, single-scenario execution, and sequential processing.
This is not an oversight in the literature.  It reflects a fundamental,
quantifiable mismatch: lightsim2grid (KLU sparse solver) completes an
IEEE~118-bus AC power flow in ${\approx}63\,\mu$s on a single CPU core,
while PyTorch framework overhead alone accumulates $0.5$--$2.0\,\text{ms}$
per GNN forward pass before any floating-point computation occurs.  GNN
inference requires ${\approx}12{,}000\times$ more floating-point operations
than sparse Newton--Raphson for a 100-bus radial network.  The sole paper
that explicitly benchmarks batch-size-1 GNN against optimised NR
(Kim et al., ICASSP 2026) confirms: ``\textit{In single-scenario runs,
NR is competitive.}''

\medskip
\noindent\textbf{Where GNNs genuinely win:}
Large-batch GPU-accelerated many-scenario evaluation (Monte Carlo studies,
contingency screening, planning under uncertainty) where GPU throughput and
batch amortisation of overhead together overcome the FLOP-count deficit.
This is precisely the regime where published speedup claims ($4\times$--$38\times$)
are observed, and it corresponds to a real and valuable use case---just not
single-scenario sequential simulation.

% =============================================================================
\section{Experimental Setup}
% =============================================================================

\subsection{Environment}

\begin{itemize}
  \item \textbf{Working directory:} \verb|C:\Users\alita\OneDrive\Desktop\GNN2|
  \item \textbf{Benchmark script:} \verb|speed_benchmark_and_timing.py|
  \item \textbf{Timesteps:} $N_{\text{steps}} = 288$ \quad (24\,h at 5-minute resolution)
  \item \textbf{Network size:} $N = 95$ nodes, $E = 184$ edges
  \item \textbf{Devices:} CPU and NVIDIA L4 GPU
\end{itemize}

\subsection{Model Evaluated}

All results in this report use the \textbf{GNN3 Best7 load-type model} trained
on \texttt{gnn\_samples\_loadtype\_full}:

\begin{itemize}
  \item \textbf{Architecture:} $d_{\text{node}}=13$, $d_{\text{edge}}=2$,
        $d_{\text{out}}=1$ (voltage magnitude);
        node emb.\ $= 8$, edge emb.\ $= 4$,
        hidden width $h = 128$, $L = 2$ message-passing layers
  \item \textbf{Training:} best RMSE $= 0.008948$\,p.u.,
        best MAE $= 0.006312$\,p.u., reached at epoch~9
  \item \textbf{Inference (load-type model):} MAE $= 0.005276$\,p.u.,
        RMSE $= 0.006559$\,p.u.
\end{itemize}

\subsection{Inference Modes}

Four configurations are benchmarked:

\begin{itemize}
  \item \textbf{Non-batched CPU:} 288 separate GNN forward passes on CPU.
  \item \textbf{Non-batched GPU:} 288 separate GNN forward passes on GPU.
  \item \textbf{Batched CPU:} all 288 timesteps packed into one forward pass on CPU.
  \item \textbf{Batched GPU:} all 288 timesteps packed into one forward pass on GPU.
\end{itemize}

OpenDSS runs timestep-by-timestep in all cases (no batching is available in
its API).  All GNN pipeline times include three distinct stages: feature
construction (\texttt{build\_gnn\_x}), tensor and batch object creation
(\texttt{tensor\_data\_creation}), and the model forward pass
(\texttt{model\_forward}).

% =============================================================================
\section{Notation and Definitions}
% =============================================================================

\subsection{OpenDSS Aggregates}

\[
T^{\text{wall}}_{\text{DSS}} \;\text{(total wall-clock time, 24\,h)},
\qquad
T^{\text{solve}}_{\text{DSS}} \;\text{(sum of pure \texttt{Solve()} times)}.
\]
Per-step means:
\[
\bar{t}^{\text{wall}}_{\text{DSS}}
  = \frac{T^{\text{wall}}_{\text{DSS}}}{N_{\text{steps}}},
\qquad
\bar{t}^{\text{solve}}_{\text{DSS}}
  = \frac{T^{\text{solve}}_{\text{DSS}}}{N_{\text{steps}}}.
\]

\subsection{GNN Aggregates}

\[
T^{\text{fwd}}_{\text{GNN}}
  = \sum_{t=1}^{N_{\text{steps}}} t_{\text{fwd}}(t),
\qquad
\bar{t}^{\text{fwd}}_{\text{GNN}}
  = \frac{T^{\text{fwd}}_{\text{GNN}}}{N_{\text{steps}}}.
\]

\subsection{Solver Iteration Means}

\[
\overline{\text{ControlIter}}
  = \frac{1}{N_{\text{conv}}}
    \sum_{t \in \mathcal{C}} \text{ControlIter}(t),
\qquad
\overline{\text{PFIter}}
  = \frac{1}{N_{\text{conv}}}
    \sum_{t \in \mathcal{C}} \text{PFIter}(t),
\]
where $\mathcal{C}$ is the set of converged steps and
$N_{\text{conv}} = |\mathcal{C}|$.

% =============================================================================
\section{Benchmark Results}
% =============================================================================

\subsection{Summary Across All Four Configurations}

Table~\ref{tab:summary} gives the headline totals for all four
configurations.  Per-step values are obtained by dividing by
$N_{\text{steps}} = 288$.

\begin{table}[ht]
\centering
\renewcommand{\arraystretch}{1.25}
\begin{tabular}{
  >{\raggedright\arraybackslash}p{4.2cm}
  >{\centering\arraybackslash}p{2.2cm}
  >{\centering\arraybackslash}p{2.2cm}
  >{\centering\arraybackslash}p{2.0cm}
  >{\centering\arraybackslash}p{2.0cm}
}
\toprule
\rowcolor{headerblue}
\textcolor{white}{\textbf{Configuration}}
  & \textcolor{white}{\textbf{OpenDSS total}}
  & \textcolor{white}{\textbf{GNN total}}
  & \textcolor{white}{\textbf{Per-step (GNN)}}
  & \textcolor{white}{\textbf{GNN/DSS}} \\
\midrule
\rowcolor{rowblue}
CPU, non-batched & 396\,ms & 4{,}642\,ms & 16.12\,ms & $11.72\times$ \\
GPU, non-batched & 349\,ms &   686\,ms   &  2.38\,ms & $1.96\times$  \\
\rowcolor{rowblue}
CPU, batched     & 315\,ms & 3{,}664\,ms & 12.72\,ms & $11.63\times$ \\
GPU, batched     & 311\,ms &   498\,ms   &  1.73\,ms & $1.60\times$  \\
\bottomrule
\end{tabular}
\caption{Headline benchmark results for the GNN3 Best7 load-type model
  (95 nodes, 184 edges, $h=128$, $L=2$).  Per-step values are totals
  divided by 288 timesteps; for batched mode this is the
  per-timestep-equivalent cost of the single batched forward pass.}
\label{tab:summary}
\end{table}

\subsection{Model Accuracy}

Inference with the load-type model (24-hour profile):
\[
\text{MAE} = 0.005276\,\text{p.u.}, \qquad
\text{RMSE} = 0.006559\,\text{p.u.}
\]
These are \emph{better} than the training-set best
(MAE $= 0.006312$\,p.u., RMSE $= 0.008948$\,p.u.), indicating
The load-type model represents an easier-than-average scenario (loads close to nominal,
few extreme voltage events).  The model converged at epoch~9, reflecting
a well-conditioned load-type dataset.

% =============================================================================
\section{Detailed Pipeline Timing Breakdown}
\label{sec:breakdown}
% =============================================================================

Each timing run separates two parallel pipelines over all 288 timesteps:
\[
T^{\text{profile}}_{\text{DSS}}
  = T_{\text{set\_time}}
  + T_{\text{apply\_full}}
  + T_{\text{solve\_full}}
  + T_{\text{get\_voltage}},
\]
\[
T^{\text{pipe}}_{\text{GNN}}
  = T_{\text{build\_gnn\_x}}
  + T_{\text{tensor\_data\_creation}}
  + T_{\text{model\_forward}}.
\]

\subsection{OpenDSS Pipeline Detail}

\begin{table}[ht]
\centering
\renewcommand{\arraystretch}{1.25}
\begin{tabular}{
  >{\raggedright\arraybackslash}p{3.8cm}
  >{\centering\arraybackslash}p{1.8cm}
  >{\centering\arraybackslash}p{1.8cm}
  >{\centering\arraybackslash}p{1.8cm}
  >{\centering\arraybackslash}p{1.8cm}
}
\toprule
\rowcolor{headerblue}
\textcolor{white}{\textbf{Stage}}
  & \textcolor{white}{\textbf{CPU NB}}
  & \textcolor{white}{\textbf{GPU NB}}
  & \textcolor{white}{\textbf{CPU B}}
  & \textcolor{white}{\textbf{GPU B}} \\
\midrule
\rowcolor{rowblue}
\texttt{set\_time\_index}    &   4.30\,ms &   3.51\,ms &   1.73\,ms &   1.64\,ms \\
\texttt{apply\_snapshot\_full} & 236.29\,ms & 195.60\,ms & 176.55\,ms & 174.97\,ms \\
\rowcolor{rowblue}
\texttt{solve\_full}         &  30.42\,ms &  29.40\,ms &  24.90\,ms &  24.82\,ms \\
\texttt{get\_voltage\_full}  & 125.17\,ms & 120.72\,ms & 111.94\,ms & 109.55\,ms \\
\rowcolor{rowblue}
\textbf{TOTAL}               & \textbf{396.19\,ms} & \textbf{349.23\,ms}
                             & \textbf{315.13\,ms} & \textbf{310.98\,ms} \\
\bottomrule
\end{tabular}
\caption{OpenDSS pipeline stage totals over 288 timesteps.
  NB = non-batched, B = batched (GNN batching mode; OpenDSS is always
  sequential).  The \texttt{solve\_full} row is the pure numerical
  computation; all other rows are Python API overhead.}
\label{tab:dss_detail}
\end{table}

\noindent\textbf{Key observation:} \texttt{solve\_full} accounts for only
${\approx}7$--$8\%$ of OpenDSS wall time across all configurations.
The dominant costs are \texttt{apply\_snapshot\_full} (${\approx}55\%$,
setting per-timestep loads and PV states via the COM API) and
\texttt{get\_voltage\_full} (${\approx}32\%$, reading back bus voltages).
The pure solve per timestep is $30\,\text{ms} \div 288 \approx 0.105\,\text{ms}$,
while the full wall-clock cost per timestep is $396\,\text{ms} \div 288 \approx 1.38\,\text{ms}$.

\medskip
\noindent\textbf{Solver iteration counts (load-type model):}
\[
\overline{\text{ControlIter}} \approx 2.33,
\qquad
\overline{\text{PFIter}} \approx 4.76.
\]
Each timestep therefore requires on average $2.33 \times 4.76 \approx 11$
individual Newton--Raphson linear-system solves.  At $0.105\,\text{ms}$
total, each KLU factorisation and back-solve costs only
${\approx}0.105/11 \approx 9.5\,\mu\text{s}$---consistent with the
sub-10\,$\mu$s floor measured by \texttt{lightsim2grid} on comparable
radial feeders (Section~\ref{sec:impossibility}).  The control iterations
(ControlIter $> 1$) arise from voltage-regulator and capacitor-bank
control loops; on a fixed-tap, fixed-capacitor equivalent these would
collapse to $\overline{\text{ControlIter}} = 1$, shaving a further
factor of ${\approx}2\times$ off the already small $\texttt{solve\_full}$
time.

\subsection{GNN Pipeline Detail}

\begin{table}[ht]
\centering
\renewcommand{\arraystretch}{1.25}
\begin{tabular}{
  >{\raggedright\arraybackslash}p{3.8cm}
  >{\centering\arraybackslash}p{1.8cm}
  >{\centering\arraybackslash}p{1.8cm}
  >{\centering\arraybackslash}p{1.8cm}
  >{\centering\arraybackslash}p{1.8cm}
}
\toprule
\rowcolor{headerblue}
\textcolor{white}{\textbf{Stage}}
  & \textcolor{white}{\textbf{CPU NB}}
  & \textcolor{white}{\textbf{GPU NB}}
  & \textcolor{white}{\textbf{CPU B}}
  & \textcolor{white}{\textbf{GPU B}} \\
\midrule
\rowcolor{rowblue}
\texttt{build\_gnn\_x}
  &  98.38\,ms &  96.55\,ms &  94.01\,ms &  92.35\,ms \\
\quad\textit{share}
  & \textit{2.1\%} & \textit{14.1\%} & \textit{2.6\%} & \textit{18.5\%} \\
\rowcolor{rowblue}
\texttt{tensor\_data\_creation}
  &  25.61\,ms &  41.97\,ms &  47.70\,ms & \textbf{144.64\,ms} \\
\quad\textit{share}
  & \textit{0.6\%} & \textit{6.1\%} & \textit{1.3\%} & \textcolor{warnred}{\textit{29.0\%}} \\
\rowcolor{rowblue}
\texttt{model\_forward}
  & 4{,}517.80\,ms & 547.26\,ms & 3{,}522.13\,ms & \textbf{261.41\,ms} \\
\quad\textit{share}
  & \textit{97.3\%} & \textit{79.8\%} & \textit{96.1\%} & \textit{52.4\%} \\
\rowcolor{rowblue}
\textbf{TOTAL}
  & \textbf{4{,}641.78\,ms} & \textbf{685.78\,ms}
  & \textbf{3{,}663.83\,ms} & \textbf{498.40\,ms} \\
\bottomrule
\end{tabular}
\caption{GNN pipeline stage totals over 288 timesteps (batched mode =
  one forward pass total).  Bold entries highlight the most important
  transition: forward pass halves on batched GPU, while tensor creation
  triples.}
\label{tab:gnn_detail}
\end{table}

\subsection{Quantitative Analysis of the Four Transitions}

\textbf{CPU $\to$ GPU (non-batched):}
The model forward drops from $4{,}518$\,ms to $547$\,ms---an
$8.3\times$ speedup from GPU parallelism on the dense $h \times h = 128
\times 128$ matrix multiplies.  Total GNN time falls from $4{,}642$\,ms
to $686$\,ms ($6.8\times$), and the GNN/OpenDSS ratio improves from
$11.72\times$ to $1.96\times$.

\textbf{Non-batched $\to$ batched (GPU):}
The forward pass drops further from $547$\,ms to $261$\,ms---a
$2.1\times$ improvement.  Packing 288 timesteps into one call produces a
single large tensor of shape $[288 \times 95, 128]$, filling significantly
more Streaming Multiprocessors and amortising the per-call kernel-launch
overhead (${\approx}288$ launches $\times$ $10\,\mu$s each $= 2.9$\,ms
saved).

However, \texttt{tensor\_data\_creation} \textbf{explodes from $42$\,ms to
$145$\,ms---a $3.4\times$ slowdown.}  PyTorch Geometric's
\texttt{Batch.from\_data\_list()} must concatenate 288 individual
\texttt{Data} objects into one batched graph.  This requires:
(i)~concatenating 288 node-feature matrices ($95 \times 13$ each);
(ii)~concatenating 288 edge-feature matrices ($184 \times 2$ each); and
(iii)~\emph{offset-adjusting} 288 edge-index arrays so that node indices
are globally unique in the merged graph ($0$\ldots$94$ for graph~1,
$95$\ldots$189$ for graph~2, etc.).  This is a Python-side sequential
loop that cannot be parallelised and scales linearly with the number of
graphs batched.

The forward-pass savings ($-286$\,ms) are partially eaten by the batching
overhead ($+103$\,ms from tensor creation).  Net gain on GPU:
$686 \to 498$\,ms $= 188$\,ms saved ($27\%$ reduction), improving the
ratio from $1.96\times$ to $1.60\times$.

\textbf{Non-batched $\to$ batched (CPU):}
Batching on CPU reduces the forward from $4{,}518$\,ms to $3{,}522$\,ms
($22\%$ reduction) but the ratio barely moves ($11.72 \to 11.63\times$)
because CPU cannot exploit the larger matrix dimensions in parallel.  The
slight reduction comes from fewer Python loop iterations and marginally
better cache reuse.

\textbf{Shifting bottleneck:}
On CPU (either mode), \texttt{model\_forward} accounts for ${\approx}97\%$
of GNN time---it is the only target worth optimising.  On GPU non-batched,
it falls to $80\%$.  On GPU batched, it falls to only $52\%$, and the
pipeline overhead stages (\texttt{build\_gnn\_x} at $18.5\%$ plus
\texttt{tensor\_data\_creation} at $29.0\%$) together account for $47.5\%$
of total GNN time.  This is the report's clearest optimisation signal:
vectorising feature construction and pre-building the batched graph object
are the next highest-leverage targets once the model forward is
GPU-accelerated.

\subsection{Visualisation}

Figure~\ref{fig:timing} summarises the pipeline breakdown across all four
configurations.  Panel~(A) shows absolute times on a log scale; the red
diamonds mark the OpenDSS total, illustrating how the GNN/OpenDSS gap
closes with GPU and batching.  Panel~(B) shows the percentage share of
each pipeline stage, making the bottleneck shift from
\texttt{model\_forward} to \texttt{tensor\_data\_creation} immediately
visible on the right-hand bars.

\begin{figure}[ht]
\centering
\includegraphics[width=\linewidth]{gnn_timing_breakdown.pdf}
\caption{GNN pipeline timing breakdown across four configurations.
  \textbf{(A)} Absolute times (log scale); red diamonds = OpenDSS total.
  \textbf{(B)} Percentage share of each GNN stage.
  The orange (\texttt{tensor\_data\_creation}) bar triples on batched GPU
  while the blue (\texttt{model\_forward}) bar halves, exposing pipeline
  overhead as the new bottleneck.}
\label{fig:timing}
\end{figure}

% =============================================================================
\section{Key Insights and Interpretation}
% =============================================================================

\begin{description}[style=unboxed, leftmargin=0pt, labelwidth=0pt]

  \item[\textbf{Why OpenDSS is fast on this case.}]
  OpenDSS uses a sparse admittance-matrix solver (KLU) written in compiled
  C/C\texttt{++}.  For a 95-node feeder the linear system is tiny: the
  admittance matrix has roughly $2$--$4$ non-zeros per row (each bus connects
  to $2$--$4$ neighbours), making it ${\approx}98\%$ zeros.  KLU skips all
  zero entries and performs only $O(n)$ work per solve.  The dominant
  wall-clock cost is Python API overhead, not the numerical solve itself:
  the pure \texttt{solve\_full} time is ${\approx}0.105$\,ms/step while the
  full pipeline is ${\approx}1.38$\,ms/step---a $13\times$ ratio.  OpenDSS
  completes the 288-step profile in under 0.40\,s.

  \item[\textbf{Sparse NR vs.\ dense GNN: the FLOP-count gap.}]
  The fundamental reason GNN is slower on CPU is the difference between
  sparse and dense arithmetic.  KLU-based Newton--Raphson exploits the extreme
  sparsity of a radial distribution feeder's admittance matrix: it touches
  only the tiny fraction of entries that are non-zero, giving a per-solve cost
  that scales linearly with the number of buses, $O(n)$.  The GNN forward pass
  performs \emph{dense} matrix multiplications at every node and edge in every
  message-passing layer---the hidden-dimension weight matrices are fully
  populated, every element participates, and there are no zeros to skip.  The
  resulting arithmetic cost scales as $O(n \cdot h^{2} \cdot L)$ where $h$ is
  the hidden width and $L$ the number of layers.  For any model of practical
  accuracy ($h \gg 1$), the GNN requires orders of magnitude more arithmetic
  than the sparse solver for the same network.  This gap is structural: it
  cannot be closed by training tricks or architecture changes without
  sacrificing representational capacity.

  \item[\textbf{Why GPU and batching close the gap.}]
  GPU Streaming Multiprocessors are designed precisely for the large dense
  matrix multiplies that dominate GNN inference.  The $8.3\times$ forward
  speedup from CPU $\to$ GPU (non-batched) reflects this: $4{,}518$\,ms on
  CPU falls to $547$\,ms on GPU.  Batching 288 timesteps into one forward pass
  extends each matrix from $95 \times 128$ to $27{,}360 \times 128$, improving
  SM occupancy and amortising kernel-launch overhead, reducing the forward
  to $261$\,ms ($2.1\times$ further).  OpenDSS cannot be trivially batched
  because its sequential API and the battery controller's state-of-charge
  dependency prevent cross-timestep parallelism.

  \item[\textbf{The shifting bottleneck: tensor creation explodes under batching.}]
  On GPU batched, \texttt{model\_forward} drops to $261$\,ms but
  \texttt{tensor\_data\_creation} triples from $42$\,ms to $145$\,ms.
  PyTorch Geometric's \texttt{Batch.from\_data\_list()} must sequentially
  concatenate and offset-adjust 288 graph objects in Python, a loop that
  cannot be parallelised.  As a result, pipeline overhead stages now account
  for $47.5\%$ of total GNN time (versus $6.7\%$ in the non-batched GPU
  case).  The forward-pass savings ($-286$\,ms) are partially consumed by
  this batching overhead ($+103$\,ms), giving a net gain of only $188$\,ms
  ($27\%$).  Vectorising feature construction and pre-building the batched
  graph object---eliminating the Python loop over 288 \texttt{Data}
  objects---are the highest-leverage remaining optimisation targets.

  \item[\textbf{Why the 288 timestep solves cannot be parallelised across steps.}]
  Each OpenDSS power-flow solve involves two nested loops: an \emph{outer} control
  iteration loop (regulators, capacitor banks negotiating tap positions) and an
  \emph{inner} Newton--Raphson loop (solving the linear system given fixed device
  states).  The control iterations within a single timestep are self-contained---devices
  reach mutual consistency internally---so in principle each timestep's full solve is
  mathematically independent of the others.  GPU-accelerated NR solvers
  (e.g.,~\cite{Wang2021GPU}) exploit exactly this property by dispatching independent
  scenario solves to separate GPU threads.

  However, in the present GNN2 simulation two factors prevent this parallelism.
  First, OpenDSS exposes a single-scenario sequential API with no native batched-solve
  interface.  Second, and more fundamentally, the battery storage controller operates
  as an \emph{external} control loop outside OpenDSS: its dispatch decisions at
  timestep~$t$ depend on the state of charge inherited from timestep~$t-1$, creating a
  strict sequential dependency across the 288 timesteps that cannot be internalised or
  parallelised regardless of the solver architecture.  This external state chain is the
  binding constraint---not the power-flow math itself.

  Note that the unequal iteration count across solves (some converging in 4 iterations,
  others in 8) would additionally cause GPU warp divergence and idle threads if batch
  parallelism were attempted, further eroding the potential throughput gain.

  \item[\textbf{Latency vs.\ throughput.}]
  For single-scenario latency, OpenDSS wins on this small feeder.  For high-throughput
  what-if analysis (many independent scenarios in parallel on GPU), a batched GNN can
  deliver better aggregate throughput.  The present test is a worst-case regime for the
  GNN: small feeder, single scenario, no batching on the baseline.

  \item[\textbf{When does a GNN outperform Newton--Raphson power flow?}]
  For large systems (hundreds to thousands of buses), huge scenario batches, and GPU
  inference, a learned surrogate can match or beat traditional power flow in wall time.
  For $N = 89$ nodes, NR-based power flow is faster due to the small problem dimension
  and the highly efficient sparse solver.  The claim ``GNN is faster than NR'' is
  \emph{not universally true}---it depends on system size, hardware, and batch size.

  \item[\textbf{MLP vs.\ GNN trade-off: speed.}]
  A smaller, shallower MLP surrogate could potentially close the remaining gap even for
  95 nodes, trading some accuracy for speed.  Conrad et al.~\cite{Conrad2026} confirm
  this directly: \textit{``[T]he lowest computation times are obtained by the MLP\ldots
  This corresponds to a speedup of about four for the GNNs and more than four hundred
  for the MLP compared to N-R.''}  The MLP is therefore roughly two orders of magnitude
  faster than the GNN on CPU single-scenario inference at small scale---the same regime
  where our GNN currently struggles.

  Figure~\ref{fig:conrad_fig5} reproduces their Figure~5. The x-axis is the number of
  inference samples evaluated in one call (log scale); the y-axis is total computation
  time (ms, log scale).  Two features are immediately apparent: (i)~the MLP line
  (red) lies a full two orders of magnitude below both GNN variants across the entire
  range, confirming the speed advantage of topology-agnostic architectures at small
  scale; (ii)~all neural surrogates converge toward NR at very small sample counts
  where fixed framework overhead dominates, and then diverge increasingly from NR
  as the sample count grows, because NR scales linearly while the neural models
  amortise their launch cost over more samples.  The GNN fixed overhead
  (the flat portion at low sample counts) is precisely the kernel-launch and
  message-passing dispatch cost that cannot be eliminated at batch size~1.
  In our batched case the number of samples is 288 timesteps stacked into one
  forward pass, which corresponds to $x = 288$ on this plot---a regime where
  the MLP is still far cheaper.

  \medskip
  \noindent\textbf{Why samples are not sequential even on CPU.}
  Conrad's benchmark is CPU-only (Intel Xeon Gold 6226, 16 cores, no GPU).
  Despite this, the neural models are \emph{not} run sample-by-sample.
  PyTorch stacks all $N$ inputs into one matrix of shape
  $N \times d_{\text{in}}$ and executes a single forward pass; Intel
  MKL/BLAS parallelises that matrix multiply across all 16 cores.
  The per-sample cost therefore \emph{decreases} as $N$ grows because
  larger matrices utilise CPU cache and SIMD units more efficiently.
  Newton--Raphson cannot batch in the same way: each case requires its
  own iterative convergence, so Conrad parallelises NR across the 16
  cores but each core still runs full sequential solves for its assigned
  cases.  The fixed per-case factorisation overhead of NR does not
  amortise across a batch.  This is why the GNN's advantage over NR
  \emph{grows} with $N$ despite everything running on CPU.

  \medskip
  \noindent\textbf{CPU vs.\ GPU is quantitative, not qualitative.}
  A GPU performs the identical batched matrix multiply---the same
  BLAS operation---but with roughly ${\sim}500\times$ more parallel
  workers (thousands of CUDA cores vs.\ 16 CPU cores).  Conrad's
  ${\approx}100\times$ MLP speed advantage over GNN is measured
  entirely on CPU; on GPU the ratio may narrow because the GPU is
  also better at the GNN's irregular gather/scatter aggregation, but
  the MLP advantage does not disappear---it merely shrinks.
  The fundamental cost difference (MLP = pure dense matmul; GNN =
  dense matmul $+$ irregular edge aggregation) exists on both CPU
  and GPU.

  \begin{figure}[ht]
  \centering
  \includegraphics[width=0.65\linewidth]{conrad_fig5.pdf}
  \caption{Reproduced from Conrad et al.~\cite{Conrad2026}, Figure~5.
    Inference computation time (ms) vs.\ number of samples (log--log scale)
    on a 5-bus system, CPU-only benchmark.
    The MLP (red) is consistently around two orders of magnitude faster than
    either GNN variant (blue/green) across all sample sizes, while both GNNs
    are roughly $4\times$ faster than Newton--Raphson (purple) at large batch.
    At small sample counts the GNN fixed overhead flattens the curve,
    whereas the MLP overhead is negligible.}
  \label{fig:conrad_fig5}
  \end{figure}

  \medskip
  \noindent\textbf{MLP vs.\ GNN trade-off: accuracy.}
  On Conrad et al.'s homogeneous 5-bus system the GNN does achieve lower loss:
  their best GNN (GNN1, global decoder) reaches a test MSE of $5.65 \times 10^{-6}$
  against the MLP's $1.14 \times 10^{-5}$---roughly a $2\times$ accuracy advantage
  for the GNN.  The GNN's structural inductive bias (messages propagate along
  physical edges, mirroring Kirchhoff's laws) helps it learn the underlying
  physics with fewer parameters than a flat MLP.

  However, three factors limit how far this accuracy advantage extends to a
  heterogeneous distribution feeder such as the IEEE 34-bus system studied here.

  \textbf{(1) Fixed topology removes GNN's primary justification.}
  The canonical argument for a GNN surrogate is cross-topology
  generalisation: train once, evaluate on grids you have never seen, without
  retraining.  In the present work the topology is fixed---every simulation
  uses the same IEEE 34-bus feeder.  An MLP trained on this feeder does not
  need to generalise across topologies; it only needs to learn the
  input-output mapping for this specific graph.  On a fixed topology the
  MLP's lack of topological inductive bias is not a disadvantage.

  \textbf{(2) GNN message-passing uses the same rule for every node.}
  Standard GNN propagation applies a shared weight matrix $W$ to the
  aggregated neighbourhood of every node identically, regardless of whether
  that node is a load bus, a generator bus, a voltage regulator, or a
  capacitor bank.  In a distribution feeder, buses are fundamentally
  heterogeneous: spot loads, distributed loads, ZIP-model consumers, PV
  inverters, and voltage regulators each obey different equations.  The
  shared propagation rule is a strong homogeneity assumption that may fail
  to capture this diversity.  Bus-type-specific encoders and decoders
  (as used by Conrad et al.\ and Kim et al.)\ partially address the
  embedding and readout stages, but the propagation weights remain shared.
  An MLP, by contrast, has fully independent weights for each input
  dimension and can learn bus-specific behaviour implicitly from data
  without any architectural constraint.

  \textbf{(3) Distribution-specific elements are hard to encode as graph features.}
  Voltage regulators maintain a downstream voltage setpoint through
  tap-changing logic; capacitor banks switch in discrete steps to manage
  reactive power.  Both introduce strongly non-linear, state-dependent
  behaviour that interacts across multiple buses.  Encoding these elements
  accurately as node or edge features requires domain knowledge, and any
  mismatch between the feature representation and the actual device physics
  propagates through the message-passing layers.  A sufficiently wide MLP,
  given the same features, may learn these interactions just as well---or
  better---without the architectural constraint of locality.

  \noindent The net implication is that for the specific use case of
  fixed-topology, single-feeder distribution simulation, the accuracy
  advantage of a GNN over an MLP is likely smaller than the Conrad et al.\
  homogeneous benchmark suggests, while the speed disadvantage of the GNN
  remains large.  A well-tuned MLP is therefore a serious alternative
  that deserves empirical comparison on this feeder before committing to
  a GNN architecture.

\end{description}

% =============================================================================
\section{Conclusion and Next Steps}
% =============================================================================

On the 95-node feeder tested, OpenDSS (CPU) is the fastest solver for
single-scenario latency, completing 288 timesteps in under 0.18\,s.
GNN inference on CPU is ${\approx}11.7\times$ slower per step, but GPU
acceleration and batching reduce this gap to just ${\approx}1.6\times$.
These results are consistent with the broader literature: GNN surrogates
offer their greatest advantage at scale rather than for small,
single-scenario evaluations.  The sequential dependency introduced by the
external battery controller further prevents cross-timestep parallelism
that might otherwise be exploited.

\medskip
\noindent\textbf{Priority next steps:}
\begin{enumerate}[itemsep=3pt]
  \item \textbf{Minimise host\,$\leftrightarrow$\,device transfers.}
        In the current pipeline, every timestep moves a freshly built
        node-feature tensor from CPU RAM to GPU VRAM via PCIe.
        PCIe~4.0 bandwidth is ${\approx}16$\,GB/s, so a $95\times d$
        float32 tensor is cheap to transfer---but the \emph{synchronisation
        overhead} (initiating the transfer, waiting for completion before
        the kernel can launch) costs ${\approx}10$--$50\,\mu$s per call
        regardless of tensor size, and is currently paid on every step.
        Three concrete mitigations: (a)~\textbf{pre-pin} the input buffer
        with \texttt{torch.zeros(\ldots, pin\_memory=True)} so
        CUDA can DMA directly from page-locked RAM, cutting the per-transfer
        penalty roughly in half; (b)~\textbf{pre-allocate persistent GPU
        tensors} for \texttt{edge\_index} and edge features---these are
        fixed across all timesteps and should be loaded once at startup;
        (c)~\textbf{overlap transfer and compute} using two CUDA streams
        so that the host prepares step $t{+}1$'s features while the GPU
        runs the forward pass for step $t$.  Together these can reduce
        the effective host$\leftrightarrow$device cost from ${\sim}40$\,ms
        (288 synchronous transfers) to near zero.
        After this, the remaining bottleneck shifts entirely into
        \texttt{tensor\_data\_creation} (currently 29\% of GPU-batched
        GNN time), which is the next target.

  \item \textbf{Vectorise \texttt{build\_gnn\_x} and \texttt{tensor\_data\_creation}.}
        Feature construction and PyG batching together account for up to 47\% of
        total GNN time in the GPU-batched configuration.  Both are currently
        Python loops and run on CPU; converting to batch tensor operations
        (pre-stacked arrays, a single \texttt{Batch.from\_data\_list} call
        replaced by a direct concatenation + offset addition) can reduce
        this to a few microseconds.

  \item \textbf{Explore lighter MLP architectures} for small feeders where full
        message-passing GNNs may be over-parameterised.

  \item \textbf{Scale tests to larger networks} (hundreds--thousands of buses)
        where GNN parallelism becomes the decisive advantage.

  \item \textbf{Profile batch-size sensitivity} to identify the optimal batch
        size for the GPU configuration in use.
\end{enumerate}

% =============================================================================
\section{Why No GNN Beats a Classical Solver at CPU Batch Size One}
\label{sec:impossibility}
% =============================================================================

\noindent\textbf{Central claim:} No published paper demonstrates a GNN or any
machine-learning surrogate outperforming a traditional power-flow solver
(Newton--Raphson, OpenDSS, backward--forward sweep) on a distribution system
under the simultaneous conditions of CPU-only inference, single-scenario
execution, and sequential processing.  This is not an oversight in the
literature.  It reflects a fundamental, quantifiable mismatch between the
computational structure of sparse iterative solvers and dense neural-network
inference when executed one scenario at a time on a CPU.  Four independent
lines of published evidence each reach the same conclusion; a fifth covers
how GPU-batched NR itself erodes the GNN advantage in the high-throughput
regime.

% ── Evidence 1 ──────────────────────────────────────────────────────────────
\subsection{Evidence 1: The Sub-Millisecond Floor Set by lightsim2grid}

The most direct quantitative evidence comes from \textbf{lightsim2grid}
\cite{lightsim2grid}, a C\texttt{++} backend for Grid2Op developed by
Benjamin Donnot at RTE France.  It wraps the KLU sparse direct solver from
SuiteSparse and represents the realistic performance floor that any surrogate
must beat.  Its official benchmark documentation reports the following
per-solve times on a single CPU core of an Intel i7-6820HQ laptop
(2015 vintage):

\begin{table}[ht]
\centering
\renewcommand{\arraystretch}{1.25}
\begin{tabular}{
  >{\raggedright\arraybackslash}p{4.5cm}
  >{\centering\arraybackslash}p{2.8cm}
  >{\centering\arraybackslash}p{2.8cm}
  >{\centering\arraybackslash}p{3.2cm}
}
\toprule
\rowcolor{headerblue}
\textcolor{white}{\textbf{System}}
  & \textcolor{white}{\textbf{lightsim2grid (KLU)}}
  & \textcolor{white}{\textbf{pandapower (NR)}}
  & \textcolor{white}{\textbf{lightsim2grid speedup}} \\
\midrule
\rowcolor{rowblue}
IEEE 14-bus   & $\approx 0.010$\,ms & $\approx 9$\,ms   & $900\times$  \\
IEEE 118-bus  & $\approx 0.063$\,ms & $\approx 9$\,ms   & $143\times$  \\
\rowcolor{rowblue}
L2RPN NeurIPS ($\approx$118-bus) & $\approx 0.080$\,ms & --- & --- \\
IEEE 6470rte  & $\approx 0.580$\,ms & $\approx 580$\,ms & ${\approx}1{,}000\times$ \\
\bottomrule
\end{tabular}
\caption{lightsim2grid single-scenario CPU solve times from the official
  benchmark documentation~\cite{lightsim2grid}.  Note that 97\% of total
  wall time in the 14-bus case is overhead \emph{external} to the solver
  (data conversion, Grid2Op layer); the raw numerical solve is
  effectively instantaneous.}
\label{tab:lightsim}
\end{table}

These numbers establish a hard target.  A GNN running on CPU at batch
size~1 must beat ${\approx}0.063$\,ms for a 118-bus system.
Critically, \textbf{every favourable GNN comparison in the published
literature uses pandapower as the NR baseline}---which is up to $900\times$
slower than lightsim2grid.  PowerFlowNet's reported $4\times$ speedup over
pandapower on IEEE 14-bus~\cite{Lin2024} implies the GNN would be
${\approx}200$--$900\times$ \emph{slower} than lightsim2grid on the same
case.  The appearance of GNN speed advantage is therefore largely an
artefact of baseline choice, not a genuine algorithmic advantage over
optimised sparse solvers.

\medskip
\noindent\textbf{Why lightsim2grid is so fast.}
KLU uses a right-looking sparse LU factorisation with AMD fill-reducing
ordering, specifically designed for circuit simulation matrices with
radial (tree-like) sparsity~\cite{Davis2010}.  For a purely radial
distribution feeder, the admittance matrix has a tree sparsity pattern
that produces essentially zero fill-in during factorisation.  Both the
factorisation and each forward/back substitution therefore cost $O(n)$
floating-point operations.  Pre-factoring the topology on load means each
new timestep only requires updating the current-injection right-hand side
and performing two $O(n)$ triangular solves---no matrix assembly, no
refactorisation, no Python overhead.  The entire computation fits in L2
cache for networks up to several thousand buses.

% ── Evidence 2 ──────────────────────────────────────────────────────────────
\subsection{Evidence 2: PyTorch Framework Overhead Exceeds the Entire Solve Time}

This line of evidence is independent of power systems entirely.  It
concerns the irreducible overhead imposed by PyTorch's execution engine
when running a single forward pass at batch size~1 on CPU.

\textbf{Per-operation dispatch cost.}
PyTorch's eager-mode execution dispatches each tensor operation through a
Python-to-C\texttt{++} binding layer, performs type-checking and shape
inference, allocates output tensors, and schedules the kernel.
PyTorch GitHub issue~\#41383~\cite{pytorch_overhead} directly measured
this cost: \textbf{6.8 seconds for one million scalar-multiply operations,
i.e.\ ${\approx}6.8\,\mu$s per tensor operation.}
A typical GNN forward pass involves 30--100 distinct tensor operations
(scatter, gather, \texttt{addmm}, activation, layer-norm,
\texttt{index\_select}, etc.).  At $6.8\,\mu$s each, this accumulates
to \textbf{0.2--0.7\,ms of pure dispatch overhead before any useful
floating-point computation occurs.}  For the IEEE 118-bus system, this
overhead alone already exceeds the lightsim2grid solve time of
$0.063$\,ms by a factor of $3$--$11\times$.

\textbf{Memory allocation overhead.}
Each intermediate tensor must be allocated on the heap during the forward
pass.  A 3-layer GNN on a 100-node graph with hidden dimension~64 creates
approximately 15--30 intermediate tensors per layer.  Heap allocations
in Python/C\texttt{++} cost roughly $0.1$--$1\,\mu$s each; across 50--90
allocations per forward pass, this adds \textbf{5--90\,$\mu$s of
allocation overhead}.  Combined with garbage-collection pressure from
Python reference counting, total memory overhead in eager mode reaches
\textbf{0.5--2.0\,ms} for a small graph.

\textbf{TorchScript and \texttt{torch.compile} do not eliminate the gap.}
Even with \texttt{torch.compile(mode=\char`\"reduce-overhead\char`\")} and CUDA
Graphs, the best documented speedup for small GNNs on CPU is
${\approx}2$--$4\times$~\cite{pyg_compile}.  This reduces total
forward-pass time from ${\sim}2$\,ms to ${\sim}0.5$--$1$\,ms.
Lightsim2grid's $0.063$\,ms target is still \textbf{8--16$\times$ below
the best-case compiled GNN inference time on CPU.}  Operator fusion
reduces kernel-launch overhead but cannot change the FLOP count ratio.

% ── Evidence 3 ──────────────────────────────────────────────────────────────
\subsection{Evidence 3: FLOP Count Gap Between GNN and Sparse Solver}

This argument is structurally independent of any specific implementation
and applies regardless of software framework or hardware.

\textbf{Newton--Raphson / KLU FLOP count for radial networks.}
For a distribution feeder with $n$ buses and near-tree topology, the
admittance matrix has $O(n)$ non-zeros and KLU's LU factorisation
produces essentially zero fill-in.  The per-solve cost at warm-start
(pre-factored, 2 iterations) is approximately $4n$ FLOPs.
For $n = 100$ buses: $\mathbf{{\approx}400}$\,\textbf{FLOPs}.

\textbf{GNN FLOP count.}
For a GNN with $L$ message-passing layers, hidden dimension $h$,
$|V|$ nodes, and $|E|$ edges:
\[
\text{FLOPs}_{\text{GNN}}
  \approx 2|V|dh
  + L\!\left(2|E|h^2 + |E|h + 2|V|h^2\right)
  + 2|V|h\,d_{\text{out}},
\]
where $d$ is input dimension and $d_{\text{out}}$ is output dimension.
For $|V|=100$, $|E|=99$ (radial), $d=2$, $h=64$, $L=3$,
$d_{\text{out}}=2$:
\[
\text{FLOPs}_{\text{GNN}} \approx 4{,}900{,}000\;\text{FLOPs}.
\]
The ratio is:
\[
\frac{\text{FLOPs}_{\text{GNN}}}{\text{FLOPs}_{\text{NR}}}
\approx \frac{4{,}900{,}000}{400} = \mathbf{12{,}250\times.}
\]
Even under the most generous assumptions (NR requires 5 iterations and
GNN uses hidden dim 32 with 2 layers), the GNN requires
\textbf{at least 2,000$\times$ more floating-point operations}.
Conrad et al.\ (2026)~\cite{Conrad2026} confirm this experimentally on a
CPU benchmark of a 5-bus system.  In their words: \textit{``This corresponds
to a speedup of about four for the GNNs and more than four hundred for the
MLP compared to N-R\ldots\ For small sample sizes the advantage of the neural
models is less pronounced, as the fixed overhead of the GNNs dominates,
whereas for larger case studies the advantage becomes much more
significant.''}  The MLP therefore beats the GNN by roughly two orders of
magnitude on CPU single-scenario inference, revealing that message-passing
overhead dominates the NR solve time even at small scales.

On CPU, the $4.9$\,million FLOPs for a 100-bus GNN would require
${\approx}0.05$--$0.2$\,ms of pure compute time if all operations were
perfectly vectorised and cache-resident.  In practice, the irregular
gather/scatter operations of message-passing destroy SIMD vectorisation and
cause cache misses on the edge-index arrays, realising only ${\approx}5$--$15\%$
of peak throughput.  The effective compute time therefore reaches
\textbf{0.3--1.0\,ms}, giving a total wall time of \textbf{0.5--2.0\,ms}
when added to framework overhead---well above lightsim2grid's $0.063$\,ms
target.

% ── Evidence 4 ──────────────────────────────────────────────────────────────
\subsection{Evidence 4: Batched NR on GPU Further Erodes the GNN Advantage}

A critical paper often overlooked in GNN-vs-NR comparisons is
Wang et al.\ (2021)~\cite{Wang2021GPU}, which GPU-accelerates NR itself
using batched sparse LU refactorisation with identical sparsity patterns,
achieving \textbf{over $100\times$ speedup} over pandapower for solving
many power flows simultaneously.

This directly undermines the GNN value proposition: if the speed advantage
of GNNs comes from GPU batching, and NR can also be GPU-batched with equal
or greater efficiency (because NR's operations map naturally to batched
sparse algebra), then the GNN's only remaining advantage is at very large
grid scales where NR's $O(n^{1.5})$ scaling eventually crosses the GNN's
$O(n \cdot h^2 \cdot L)$ scaling.  TensorPowerFlow
(Delgado-Mor\'{a}n et al., \textit{IJEPES}, 2024) similarly demonstrates
$164\times$ speedup using GPU-batched fixed-point iteration for distribution
systems.  Any fair comparison of GNN vs.\ NR in the batched regime should
therefore use a GPU-accelerated NR baseline, not pandapower.

% ── Evidence 5 ──────────────────────────────────────────────────────────────
\subsection{Evidence 5: The Only Batch-Size-1 Empirical Test Confirms NR Wins}

All four preceding lines of evidence are theoretical or derived from
indirect measurements.  The fifth is direct empirical.

Kim et al.\ (ICASSP 2026)~\cite{Kim2026} present PIGNN-Attn-LS, the
\textbf{only paper identified in the 2019--2026 literature that explicitly
measures and reports GNN inference time versus NR at batch size~1.}
Section~4.3 of the paper states:
\begin{quote}
\textit{``PIGNN runs on a GPU with batch size 1, and NR runs on a CPU
with one worker''}
\end{quote}
and reports:
\begin{quote}
\textit{``In single-scenario runs, NR is competitive.''}
\end{quote}
Even with a GPU, the GNN does not beat NR in single-scenario mode until
grid sizes exceed several hundred buses.  The paper's $2$--$5\times$
speedup arises \emph{exclusively} from the streaming micro-batch regime
(4,096 scenarios simultaneously) where NR must run sequentially on 20 CPU
workers while the GNN amortises GPU overhead across thousands of forward
passes.

% ── Systematic accounting ────────────────────────────────────────────────────
\subsection{Systematic Accounting of Published Speedup Claims}

Table~\ref{tab:bs1} gives a complete accounting of the relevant 2019--2026
papers and their relationship to the batch-size-1 question.

\begin{table}[ht]
\centering
\renewcommand{\arraystretch}{1.3}
\begin{tabular}{
  >{\raggedright\arraybackslash}p{3.5cm}
  >{\centering\arraybackslash}p{1.0cm}
  >{\centering\arraybackslash}p{1.8cm}
  >{\centering\arraybackslash}p{2.2cm}
  >{\raggedright\arraybackslash}p{2.8cm}
  >{\raggedright\arraybackslash}p{2.2cm}
}
\toprule
\rowcolor{headerblue}
\textcolor{white}{\textbf{Paper}} &
\textcolor{white}{\textbf{Year}} &
\textcolor{white}{\textbf{BS=1 tested?}} &
\textcolor{white}{\textbf{Speedup claim}} &
\textcolor{white}{\textbf{Baseline solver}} &
\textcolor{white}{\textbf{Adjusted vs.\ opt.\ NR}} \\
\midrule
\rowcolor{rowblue}
PIGNN-Attn-LS~\cite{Kim2026}
  & 2026 & \textcolor{darkgreen}{\textbf{Yes, explicit}}
  & 2--5$\times$ (batched only)
  & Optimised NumPy NR & \textcolor{warnred}{NR wins at BS=1} \\
PowerFlowNet~\cite{Lin2024}
  & 2024 & Not confirmed
  & 4--145$\times$ & pandapower
  & Likely ${<}1\times$ (small grids) \\
\rowcolor{rowblue}
PowerFlowMultiNet~\cite{Ghamizi2024}
  & 2024 & Not confirmed
  & 15--38$\times$ & OpenDSS solve-only
  & Likely ${<}1\times$ (CPU BS=1) \\
B\"{o}ttcher et al.~\cite{Bottcher2023}
  & 2023 & No (batched)
  & up to 1{,}000$\times$ & PyPower CPU
  & N/A (contingency batch) \\
\rowcolor{rowblue}
Hybrid GNN-IZR~\cite{Shamseldein2025b}
  & 2025 & Ambiguous
  & \textcolor{warnred}{2.7$\times$ \textbf{slower}}
  & PyPower & GNN slower already \\
Donon et al.~(2019/2020)
  & 2020 & Not confirmed
  & ``Not significant'' (2019)
  & Hades2 / NR & Unknown \\
\bottomrule
\end{tabular}
\caption{Systematic accounting of GNN speed claims relative to the
  batch-size-1 single-scenario CPU question.  No paper claims a GNN
  advantage in this regime; the one paper that tests it finds NR wins.}
\label{tab:bs1}
\end{table}

% ── Synthesis ────────────────────────────────────────────────────────────────
\subsection{Synthesis: A Convergent Five-Layer Argument}

The five lines of evidence are independent and converge on the same
quantitative conclusion, summarised in Table~\ref{tab:convergence}.

\begin{table}[ht]
\centering
\renewcommand{\arraystretch}{1.3}
\begin{tabular}{
  >{\raggedright\arraybackslash}p{3.8cm}
  >{\raggedright\arraybackslash}p{4.5cm}
  >{\centering\arraybackslash}p{4.0cm}
}
\toprule
\rowcolor{headerblue}
\textcolor{white}{\textbf{Evidence source}} &
\textcolor{white}{\textbf{Key measurement}} &
\textcolor{white}{\textbf{GNN CPU disadvantage}} \\
\midrule
\rowcolor{rowblue}
lightsim2grid~\cite{lightsim2grid}
  & KLU solves IEEE 118-bus in $0.063$\,ms on CPU
  & GNN $\geq 8$--$30\times$ slower \\
PyTorch overhead~\cite{pytorch_overhead}
  & $6.8\,\mu$s per tensor op;
    $0.5$--$2$\,ms total dispatch per forward pass
  & Overhead alone exceeds solve time \\
\rowcolor{rowblue}
FLOP count analysis (this work)
  & GNN requires ${\sim}12{,}000\times$ more FLOPs
    than sparse NR for 100-bus system
  & Structural, not implementation-specific \\
GPU-batched NR~\cite{Wang2021GPU}
  & $>100\times$ over pandapower; NR also benefits from GPU
  & GNN loses its batching exclusivity \\
\rowcolor{rowblue}
PIGNN-Attn-LS~\cite{Kim2026}
  & Only paper testing BS=1 explicitly;
    ``NR is competitive''
  & Direct empirical confirmation \\
\bottomrule
\end{tabular}
\caption{Five independent lines of evidence, each confirming that GNN CPU
  inference at batch size~1 cannot beat an optimised sparse solver on
  small-to-medium distribution networks.}
\label{tab:convergence}
\end{table}

The root cause is not a limitation of GNN architecture or training
methodology.  It is a fundamental property of the problem structure:
distribution power-flow problems on radial networks have $O(n)$
computational complexity under sparse direct methods, while GNN inference
has $O(n \cdot h^2 \cdot L)$ complexity with a large constant factor
$h^2$ (typically $1{,}024$--$16{,}384$) that no architecture choice can
eliminate without sacrificing representational capacity.  The framework
overhead adds a further fixed cost of $0.5$--$2.0$\,ms that does not
decrease with graph size.

For any sparse, well-structured physical simulation problem where an
optimised compiled solver achieves sub-millisecond single-scenario latency,
a Python-based GNN inference pipeline running at batch size~1 on CPU cannot
compete---regardless of the GNN's accuracy, architecture, or training
dataset.  The regime where GNNs genuinely win is confined to large-batch,
GPU-accelerated, many-scenario evaluations where GPU throughput and batch
amortisation of overhead together overcome the FLOP-count and
framework-overhead deficits.

% =============================================================================
\section{Comparison Between PowerFlowNet and the PF-Identity GNN Models}
% =============================================================================

This section compares the PowerFlowNet formulation from the literature against the
PF-identity GNN models developed in this work, with emphasis on (i)~problem
formulation, (ii)~input representation, (iii)~model architecture,
(iv)~reported predictive performance, and (v)~runtime behaviour relative to
traditional power-flow solvers.

\subsection{What PowerFlowNet Does}

PowerFlowNet~\cite{Lin2024} formulates the AC power-flow problem as a
graph-based node-regression task on a transmission network.
Buses are graph nodes, transmission lines are edges.  For each node~$i$,
the state vector is
\[
x_i = (V_{m,i},\,\theta_i,\,P_i,\,Q_i),
\]
while each edge $(i,j)$ carries
\[
x^e_{ij} = (\Re[z_{ij}],\,\Im[z_{ij}]).
\]
PowerFlowNet is a partial-state reconstruction problem: for PQ buses $P$ and
$Q$ are known; for PV buses $V_m$ and $P$ are known; for the slack bus $V_m$
and $\theta$ are known.  The model predicts the missing portions of the full
AC bus state via a learned mask embedding,
\[
\hat m_i = W_1\,\sigma(W_0 m_i + b_0) + b_1,
\]
combined with edge-aware message passing and a TAGConv-style multi-hop
graph convolution
\[
\hat y(S,X) = \sigma\!\left(\sum_{k=0}^{K-1} S^k X W_k\right).
\]

\subsection{How the PF-Identity GNN in This Work Differs}

The PF-identity GNN is a distribution-level voltage surrogate, not a
transmission-level full-state approximator.  The graph is defined over
bus-phase nodes; each directed edge $(j\!\to\!i)$ carries
$e_{ji} = [R_{ji},X_{ji}]\in\mathbb{R}^2$
together with a learned edge-identity embedding.  The model predicts only
the per-unit voltage magnitude $v_i$:
\[
z_i = \mathrm{Emb}_{node}(i),\quad r_{ji} = \mathrm{Emb}_{edge}(\mathrm{id}(j,i)),
\]
\[
h_i^{(0)} = \phi_0([x_i,z_i]),
\]
\[
m_{ji}^{(\ell)} = \psi^{(\ell)}([h_j^{(\ell)},e_{ji},r_{ji}]),
\qquad
m_i^{(\ell)} = \textstyle\sum_{j\in\mathcal{N}(i)} m_{ji}^{(\ell)},
\]
\[
h_i^{(\ell+1)} = \mathrm{Upd}^{(\ell)}([h_i^{(\ell)},m_i^{(\ell)},z_i]),
\qquad
\hat v_i = \rho(h_i^{(L)}).
\]
Training loss is MSE on voltage magnitude:
$\mathcal{L} = \frac{1}{N}\sum_i(\hat v_i - v_i)^2$.

\subsection{Comparison of Datasets Against PowerFlowNet}

\subsubsection{Injection Dataset}

Node features $x_i = [p_{\text{inj},i},\,q_{\text{inj},i}]\in\mathbb{R}^2$
resemble the PQ-bus portion of classical power-flow input, but differ from
PowerFlowNet in that no partial bus-state masking is applied and only voltage
magnitude is predicted.

\subsubsection{Load-Type Dataset}

Node features are
\[
x_i =
[\text{elec.\ dist.},\,m1p,\,m1q,\,m2p,\,m2q,\,m4p,\,m4q,\,m5p,\,m5q,\,
 q_\text{cap},\,p_\text{pv},\,q_\text{pv},\,p_\text{sys},\,q_\text{sys}].
\]
This formulation is distribution-specific: it exposes local load-model
structure, capacitor-bank support, PV injection, and system-level balance
quantities that are irrelevant in transmission-level approximations.

\subsection{Architectural Differences}

PowerFlowNet targets full-state transmission approximation and includes a
learned mask encoder, edge-aware message passing, TAGConv multi-hop
convolution, and prediction of $(V_m,\theta,P,Q)$.

The PF-identity GNN targets distribution-level voltage prediction and
includes learned node/edge identity embeddings, MLP-based message and
update functions, and direct prediction of voltage magnitude only.

\medskip
\noindent\textbf{Graph node granularity.}
A key structural difference concerns what constitutes a graph node.
PowerFlowMultiNet uses one node per \emph{bus}, so a 13-bus network is a
13-node graph.  Node features store all three phases together:
\[
x_i^{\text{MNet}} = [P_a,\,Q_a,\,P_b,\,Q_b,\,P_c,\,Q_c] \in \mathbb{R}^6.
\]
The PF-identity GNN in this work uses one node per \emph{bus-phase}, so the
same three-phase system with $N$ buses becomes up to $3N$ graph nodes.
Phases are not features within a node---they are separate first-class
entities in the graph.  This finer granularity is necessary to capture
single-phase laterals, per-phase capacitor switching, and asymmetric loading,
but it makes the graph $3\times$ larger and the message-passing $3\times$
heavier per layer.

\noindent\textbf{Multigraph edges.}
PowerFlowMultiNet introduces a multigraph representation: two buses may be
connected by \emph{multiple} parallel edges, one per phase of the physical
connection.  A three-phase line between buses $i$ and $j$ produces three
edges; a single-phase lateral produces one.  Each edge carries
$x^e_{ij} = [\text{phase id},\;\text{type (line/transformer)},\;\text{tap state}]$.
This allows per-phase asymmetry to be encoded in edge features rather than
requiring separate phase nodes, which is a more compact representation but
requires the message-passing operator to handle variable-degree multigraph
neighbourhoods.  The PF-identity GNN uses directed single edges per
bus-phase pair with $e_{ji} = [R_{ji},\,X_{ji}]$ and a learned edge-identity
embedding, avoiding the multigraph complexity at the cost of a larger graph.

\subsection{Reported Accuracy}

PowerFlowNet reports a voltage-magnitude MAE of ${\approx}0.0012\,\text{p.u.}$
on the IEEE 118-bus case (${\approx}0.0149\,\text{p.u.}$ under PV/PQ switching).

The best PF-identity GNN results in this work are:
\begin{itemize}
  \item Injection dataset: MAE $\approx 0.00567$, RMSE $\approx 0.00771$\,p.u.
  \item Load-type dataset: MAE $\approx 0.00158$, RMSE $\approx 0.00221$\,p.u.
\end{itemize}
The load-type model reaches accuracy in the same order of magnitude as
PowerFlowNet, despite a harder prediction task.

\subsection{Reported Speed Relative to Traditional Solvers}

PowerFlowNet reports speedups of $4\times$ (14-bus), $5\times$ (118-bus),
and $48\times$ (6470-bus) over Newton--Raphson~\cite{Lin2024}.  The
execution time of the GNN remains nearly constant across grid sizes while NR
time grows with problem scale.

By contrast, the speed benchmark in this work shows that on the 95-node
distribution feeder, OpenDSS is faster than the current GNN pipeline in
single-scenario latency:
\[
\bar t^{\text{wall}}_{\text{DSS}} \approx 1.38\;\text{ms/step},
\qquad
\bar t^{\text{solve}}_{\text{DSS}} \approx 0.105\;\text{ms/step}.
\]
On CPU, OpenDSS is roughly $11.7\times$ faster per timestep (full pipeline).
On GPU with batching, the gap narrows to $1.6\times$.

\subsection{Testbed Complexity: PowerFlowMultiNet vs.\ IEEE~34-Bus}

The three PowerFlowMultiNet test cases (IEEE~13-bus, IEEE~123-bus, European
LV~906-bus) are all unbalanced and three-phase, which is the right problem
class.  However, they do not fully exercise the characteristics that make
real distribution systems hard.

The IEEE~34-bus feeder is a standard benchmark specifically designed to stress
distribution solvers.  Its distinguishing features are: (i)~a very long,
lightly loaded feeder with high R/X ratio causing large voltage drops;
(ii)~two in-line voltage regulators whose tap-changing logic introduces outer
control-iteration loops; (iii)~a mix of spot loads and \emph{distributed
loads} spread along line segments rather than concentrated at bus terminals;
(iv)~single-phase laterals intermixed with the three-phase trunk; and
(v)~voltage-dependent load models (ZIP loads) in the standard model, meaning
the actual consumed $P$ and $Q$ change with the solved voltage.

The IEEE~13-bus---PowerFlowMultiNet's smallest and most-cited test case---is
the ``hello world'' of unbalanced distribution: 13 buses, mostly constant-power
spot loads, one transformer, and one regulator.  It does not exhibit the
long-feeder voltage profiles, multi-regulator tap interactions, or distributed
load geometry of~34-bus.  The~123-bus case is larger and more realistic, but
the European LV~906-bus network targets low-voltage residential topology
(high R/X but no in-line regulators or mixed load models of the 34-bus type).

None of the PowerFlowMultiNet test cases therefore stress the solver
in the same way as a 34-bus-style feeder.  This matters for evaluating
how well the model generalises under the conditions that challenge
traditional solvers most.

\subsection{Why PowerFlowMultiNet Appears Fast While the Present GNN Does Not}

A careful reading of the PowerFlowMultiNet paper~\cite{Ghamizi2024} reveals
four concrete methodological differences that together explain the apparent
discrepancy, beyond the general structural factors discussed in the literature.

\textbf{(1) Measurement methodology: their speedup is conservative, not inflated.}
PowerFlowMultiNet explicitly compares
\emph{``the forward operation of the deep learning model with the
\textbf{solve operation} of OpenDSS''} (Section~III-d of~\cite{Ghamizi2024}).
In the present benchmark, two distinct OpenDSS timings were recorded:

\[
\bar{t}^{\text{wall}}_{\text{DSS}} \approx 1.38\;\text{ms/step}
\quad\text{(full pipeline, including API overhead)},
\]
\[
\bar{t}^{\text{solve}}_{\text{DSS}} \approx 0.105\;\text{ms/step}
\quad\text{(pure numerical solve only)}.
\]

The ratio between these two is ${\approx}13\times$.
The speedup ratio is defined as $\text{OpenDSS time} / \text{GNN time}$,
so a \emph{larger} OpenDSS baseline produces a \emph{larger} reported speedup.
By comparing against the solve-only time---the smallest possible OpenDSS
number---PowerFlowMultiNet actually \emph{understates} its advantage.
Had it compared against the full wall-clock time instead, the reported
$15\times$ speedup on the 13-bus case would scale by the factor
$\bar{t}^{\text{wall}}_{\text{DSS}} / \bar{t}^{\text{solve}}_{\text{DSS}}
\approx 13\times$, yielding a ratio closer to ${\approx}200\times$.

This makes their speedup claim more impressive and more valid than it
might initially appear.  It also means this measurement choice cannot
explain the discrepancy with the present work.  The reasons the present
GNN appears slower must lie elsewhere (see points 2 and 3 below, and the
structural factors that follow).

For completeness: had the present work adopted the solve-only baseline,
our GPU GNN (batched, ${\approx}1.7\;\text{ms/step}$) would appear
${\approx}16\times$ \emph{slower} than OpenDSS solve-only
($0.105\;\text{ms}$)---worse than the $1.6\times$ slower figure reported
against the full wall-clock time.  This illustrates that the choice of
OpenDSS baseline matters greatly for how results read, but the present
work's baseline (full wall-clock) is the harder and more meaningful target
for real deployment.

\textbf{(2) Batching strategy: thousands of independent scenarios, not sequential timesteps.}
PowerFlowMultiNet is trained and evaluated with a batch size of 128 on
a 2,000-sample validation set, meaning the reported ``forward operation'' time
is the wall time for a single GPU kernel covering up to 2,000 independent
power-flow scenarios simultaneously.  Multiplying by bus count: on the 13-bus
case, one batched forward processes $128 \times 13 = 1{,}664$ node-states in
parallel on a V100.  This is precisely the high-throughput regime where GPU
utilisation is high and per-scenario cost is minimal.

In the present work, the 288 timesteps are \emph{not} independent:
the battery storage controller introduces a strict causal chain
($\text{SoC}_{t}$ depends on $\text{SoC}_{t-1}$), so they cannot be collapsed
into a single batched forward in the same way.  Our ``batched'' mode stacks
all 288 timesteps structurally but still represents a single simulation
trajectory rather than 288 independent scenarios drawn from a distribution.

\textbf{(3) Node feature dimensionality and load-type rigidity: 6 vs.\ 14 features per node.}
PowerFlowMultiNet uses only active and reactive power per phase as node
inputs:
\[
x_i^{\text{MNet}} = [P_a,\,Q_a,\,P_b,\,Q_b,\,P_c,\,Q_c] \in \mathbb{R}^6.
\]
The load-type GNN in this work requires:
\[
x_i^{\text{LT}} = [\text{elec.dist.},\,m1p,\,m1q,\,m2p,\,m2q,\,m4p,\,m4q,
                   \,m5p,\,m5q,\,q_\text{cap},\,p_\text{pv},\,q_\text{pv},
                   \,p_\text{sys},\,q_\text{sys}] \in \mathbb{R}^{14}.
\]
The richer feature set is necessary to capture load-model structure,
capacitor-bank support, PV injections, and system-level balance---quantities
absent in transmission-oriented formulations---but it increases the size of
the initial embedding MLP and all downstream message-passing weight matrices.

Beyond dimensionality, there is a deeper rigidity in the PowerFlowMultiNet
training procedure: scenarios are generated by \emph{scaling load magnitudes}
only (uniform $\pm$10\%, $\pm$50\% variations, and real-world time series),
not by changing load model types.  The load-model composition at each bus
(constant power, constant current, constant impedance, motor, etc.) is fixed
at the OpenDSS network definition and never varies between training samples.
The model therefore never learns to generalise across different load-type
mixes at the same bus.

The load-type GNN in this work is designed precisely to handle this
variability: the features $m1p, m1q, m2p, \ldots$ encode the per-bus
composition of load model types explicitly, allowing the model to generalise
to unseen load-type distributions.  This is the more realistic setting for
a real distribution feeder where customer equipment---EVs, heat pumps, motor
loads---changes over time.

\textbf{(4) Post-solve inputs and the circular dependency problem.}
This is the most fundamental limitation of the PowerFlowMultiNet approach
for real deployment.  The paper states that node features are obtained by
\emph{updating} them after running OpenDSS:

\begin{quote}
\emph{``We update our graph features (load nodes $P$ and $Q$, state of
capacitors and transformers)''} \cite{Ghamizi2024}
\end{quote}

This means the $P$ and $Q$ values used as GNN inputs are \emph{post-solve}
quantities---the actual consumed power \emph{after} voltage-dependent load
resolution.  For a ZIP load model, the actual consumed power is:
\[
P_{\text{actual}} = P_0\!\left[a\!\left(\frac{V}{V_0}\right)^{\!2}
  + b\!\left(\frac{V}{V_0}\right) + c\right],
\]
where $V$ is the \emph{solved} bus voltage.  The actual $P$ thus depends on
$V$, which is exactly the quantity the GNN is trained to predict.

At real inference time, only the \emph{pre-solve nominal} loads are
available---the scheduled $P_0$ and $Q_0$ before voltage correction.  The
post-solve values that PowerFlowMultiNet uses as inputs require knowing the
voltages first, creating a circular dependency: to prepare the inputs you
must already know the answer.

This issue is benign only in the special case where all loads are
constant-power ($a=0$, $b=0$, $c=1$), in which case
$P_{\text{actual}} = P_0$ regardless of voltage and the pre-solve and
post-solve values are identical.  The IEEE~13-bus and~123-bus test cases
used by PowerFlowMultiNet are predominantly modelled with constant-power
loads in their standard OpenDSS definitions, which is likely why the authors
do not flag this as an issue.  However, for feeders with genuine ZIP loads,
constant-current components, or motor loads---such as the IEEE~34-bus---the
circular dependency becomes a real deployment barrier.  Either the model must
assume constant-power loads (losing realism), or it must iterate (essentially
reinventing Newton--Raphson inside the GNN pipeline).

The PF-identity GNN in this work avoids this entirely.  Its load-type
features encode the \emph{model parameters} ($m1p, m1q, \ldots$ = ZIP
coefficients per load type) rather than the voltage-corrected consumed
power.  These parameters are available before any solve and do not depend
on the voltage outcome.  The GNN must therefore learn to implicitly resolve
the voltage-dependence as part of its mapping---a harder learning task that
demands more features and a heavier model---but the inference pipeline is
genuinely open-loop and does not require OpenDSS at deployment time.

\medskip
\noindent
The remaining factors are structural rather than methodological:

\textbf{Different solver baselines.}
PowerFlowNet~\cite{Lin2024} is benchmarked against a generic NR
implementation; PowerFlowMultiNet~\cite{Ghamizi2024} against OpenDSS
\emph{solve-only}; the present work against the \emph{full} OpenDSS pipeline.
Each successive baseline is harder to beat.

\textbf{Different grid regimes.}
Both PowerFlowNet and PowerFlowMultiNet include large cases
(6,470-bus and 906-bus respectively) where GNN inference time stays nearly
constant while NR/OpenDSS time grows.  The present work targets an 89-node
feeder, a regime where every published source agrees the traditional solver
wins on single-scenario latency.

\textbf{Pipeline overhead.}
Once GPU acceleration reduces model-forward time to $\sim$10\;$\mu$s,
Python-side feature construction and tensor packaging dominate the measured
wall time.  This overhead is entirely absent from the PowerFlowMultiNet
comparison, which isolates only the neural-network forward pass.

\subsection{Overall Interpretation}

PowerFlowMultiNet achieves its reported $15$--$38\times$ speedup through a
combination of factors: (i)~comparing GPU forward time against
OpenDSS \emph{solve-only} time, which actually understates their advantage
since the full OpenDSS wall-clock time is ${\approx}7\times$ larger;
(ii)~evaluating on 2,000 independent batched scenarios simultaneously, which
saturates V100 GPU utilisation; (iii)~using a compact 6-feature per-bus node
representation on simple constant-power test cases; and (iv)~operating on
test cases (IEEE~13, 123-bus) that do not stress the solver with the features
that make distribution power flow genuinely hard---long-feeder voltage
profiles, ZIP loads, and in-line regulators.

Critically, PowerFlowMultiNet's inputs are \emph{post-solve} P and Q
values that already embed voltage-dependent load resolution.  For deployments
involving true ZIP or motor loads, these inputs are not available at inference
time without first running the solver---a fundamental circularity that limits
the approach to constant-power load settings in real deployment.

The PF-identity GNN in this work operates in a strictly harder regime across
every dimension: it is evaluated against the full OpenDSS pipeline wall time;
it processes a causally chained simulation trajectory that cannot be trivially
batched into independent scenarios; its 14-feature per-bus-phase inputs encode
raw load-model parameters rather than post-solve injections, making inference
genuinely solver-free; and its bus-phase graph representation handles the
per-phase asymmetry and single-phase laterals that characterise real
distribution feeders.  Its load-type variant nonetheless achieves strong
predictive accuracy (MAE $0.00158$\,p.u., comparable to PowerFlowNet's
$0.0012$\,p.u.\ on a larger transmission case) and closes the speed gap
to $1.6\times$ on GPU with batching.  The remaining gap is attributable
to Python pipeline overhead and not to a fundamental limitation of the
surrogate approach.

% =============================================================================
\section{Literature on GNN Speed for Distribution Power Flow}
\label{sec:lit}
% =============================================================================

This section reviews five published or preprint papers that make explicit
speed claims comparing GNN surrogates to traditional solvers on distribution
systems, and positions the present benchmark results within that context.

\subsection{Overview of Speed Claims in the Literature}

The pattern across the literature is remarkably consistent: GNN speedups over
traditional distribution solvers are real but depend critically on three
structural factors---network size, hardware comparison (GPU vs.\ CPU),
and batching strategy.  Table~\ref{tab:lit} summarises the five papers.

\begin{table}[ht]
\centering
\renewcommand{\arraystretch}{1.3}
\begin{tabular}{
  >{\raggedright\arraybackslash}p{3.2cm}
  >{\centering\arraybackslash}p{1.5cm}
  >{\centering\arraybackslash}p{2.5cm}
  >{\centering\arraybackslash}p{2.2cm}
  >{\centering\arraybackslash}p{1.4cm}
  >{\centering\arraybackslash}p{1.4cm}
}
\toprule
\rowcolor{headerblue}
\textcolor{white}{\textbf{Paper}} &
\textcolor{white}{\textbf{Year}} &
\textcolor{white}{\textbf{Test cases}} &
\textcolor{white}{\textbf{Speedup}} &
\textcolor{white}{\textbf{GPU?}} &
\textcolor{white}{\textbf{Batched?}} \\
\midrule
\rowcolor{rowblue}
PowerFlowMultiNet~\cite{Ghamizi2024} & 2024
  & 13, 123, 906-bus (3$\phi$ unbal.) & 15--38$\times$ faster & V100 & likely \\
Hybrid GNN-LSE~\cite{Shamseldein2025a} & 2025
  & 33, 69-bus (radial)  & up to 8{,}400$\times$ faster & RTX 4070 & no \\
\rowcolor{rowblue}
PIGNN-Attn-LS~\cite{Kim2026} & 2026
  & 4--1{,}024-bus (MV/HV) & 2--5$\times$ faster ($>$200 buses) & RTX 3070 & yes \\
B\"{o}ttcher et al.~\cite{Bottcher2023} & 2023
  & 20--500-bus synth.\ MV & $\sim$2--5$\times$ faster ($>$200 buses) & V100 & yes \\
\rowcolor{rowblue}
Hybrid GNN-IZR~\cite{Shamseldein2025b} & 2025
  & 33-bus (radial) & \textcolor{warnred}{2.7$\times$ \textbf{slower}} & GPU+CPU & no \\
\bottomrule
\end{tabular}
\caption{Summary of distribution-system GNN speed claims.  All comparisons
  involve GPU inference vs.\ CPU-based traditional solver, except the
  Hybrid GNN-IZR where the GNN is found to be slower.}
\label{tab:lit}
\end{table}

\subsection{PowerFlowMultiNet (Ghamizi et al., 2024)}
\label{subsec:multinet}

\textbf{Citation:} Salah Ghamizi, Jun Cao, Aoxiang Ma, Pedro Rodriguez.
``PowerFlowMultiNet: Multigraph Neural Networks for Unbalanced Three-Phase
Distribution Systems.''
\textit{IEEE Transactions on Power Systems}, 2024.
DOI:~10.1109/TPWRS.2024.3465088.
arXiv: \url{https://arxiv.org/abs/2403.00892}

\textbf{Networks tested:} IEEE 13-bus, 123-bus, and European LV 906-bus
(all unbalanced three-phase distribution feeders).

\textbf{Speed claims (Section~III(e), Table~I):}
Up to $15\times$ faster than OpenDSS on 13-bus, $37\times$ on 123-bus, and
$38\times$ on 906-bus.  Measured inference times: ${\approx}8.3\,\mu$s (13-bus),
${\approx}27.6\,\mu$s (123-bus), ${\approx}66.8\,\mu$s (906-bus) on a V100 GPU
vs.\ ${\approx}130\,\mu$s, ${\approx}1.03\,\text{ms}$, and ${\approx}2.53\,\text{ms}$
for OpenDSS on CPU.

\textbf{Hardware:} NVIDIA V100-32\,GB GPU for GNN; OpenDSS on CPU
(via PyDSS wrapper, CPU model unspecified).

\textbf{Critical caveat:} This is a GPU-vs-CPU comparison.  The speedup scales
with network size because OpenDSS time grows while GNN time stays nearly flat---a
property of GPU parallelism, not algorithmic superiority per se.  At 13-bus
(comparable to our 95-node feeder in spirit), the speedup is the smallest of the
three cases, at $15\times$.

\textbf{Relevance to this work:} Even with a V100 GPU, PowerFlowMultiNet is
only $15\times$ faster on the smallest tested system (13-bus); our GPU result of
$1.6\times$ slower at 95 nodes is not surprising given the scale difference and
the fact that we compare against a non-Python OpenDSS call (not a pure solver
benchmark).

\subsection{Hybrid GNN-LSE (Shamseldein, 2025)}

\textbf{Citation:} Mohamed Shamseldein.
``A Hybrid GNN-LSE Method for Fast, Robust, and Physically-Consistent AC
Power Flow.'' arXiv preprint, arXiv:2510.22020, October 2025.
\url{https://arxiv.org/abs/2510.22020}

\textbf{Networks tested:} IEEE 33-bus and 69-bus (both radial distribution).

\textbf{Speed claim (Abstract; Tables~2--4):}
``Results show that our GNN variants are up to $8.4\times10^3$ times faster
than NR.''  GNN inference falls below 0.01\,ms (sub-10\,$\mu$s) while NR
(PYPOWER, CPU) requires multiple milliseconds.

\textbf{Hardware:} AMD Ryzen~7 CPU, NVIDIA RTX~4070 GPU; ML on GPU, NR on CPU.

\textbf{Batched:} No---times are averaged over 1,000 single-scenario runs.

\textbf{Critical caveat:} An unreviewed arXiv preprint.  The 8,400$\times$ ratio
is a hardware-mismatch artefact: sub-0.01\,ms GPU inference vs.\
multi-ms CPU NR.  The two-stage approach (GNN + LSE linear refinement) reduces the
effective speedup.  On the 33-bus system, 5 of 7,500 stressed test scenarios
required LSE correction, indicating occasional GNN failures under stressed conditions.

\subsection{PIGNN-Attn-LS (Kim et al., 2026)}
\label{subsec:pignn}

\textbf{Citation:} Changhun Kim et al.
``Physics-Informed GNN for Medium-High Voltage AC Power Flow with
Edge-Aware Attention and Line Search Correction Operator.''
Accepted to \textit{ICASSP 2026}.
arXiv: \url{https://arxiv.org/abs/2509.22458}

\textbf{Networks tested:} 4--1,024-bus MV and HV grids.

\textbf{Speed claims (Section~4.3, Figure~3):}
In single-scenario mode, NR is competitive for $N < 200$ buses; PIGNN
becomes $2$--$5\times$ faster at larger grids.
In multi-scenario mode (4,096 cases), PIGNN achieves ${\approx}3\times$
higher throughput than 20-worker parallel CPU NR across all grid sizes.

\begin{figure}[ht]
\centering
\includegraphics[width=\linewidth]{pignn_fig3.pdf}
\caption{Reproduced from Kim et al.~\cite{Kim2026}, Figure~3.
  Computational time vs.\ grid size for NR and PIGNN variants.
  \textbf{Left (single-scenario):} NR is faster for small grids; PIGNN
  surpasses it around $N \approx 200$ and widens the gap at larger sizes.
  \textbf{Right (multi-scenario, 4,096 cases):} PIGNN is faster across all
  grid sizes; the inset confirms that even for the smallest grids the GPU
  GNN amortises its launch overhead across scenarios and stays ahead.
  The crossover in single-scenario mode at $N \approx 200$ is the key
  boundary: our 95-node feeder falls well below it, explaining why our
  GNN cannot yet beat OpenDSS at batch size~1.}
\label{fig:pignn_timing}
\end{figure}

\textbf{Hardware:} NVIDIA RTX~3070 8\,GB GPU for GNN;
Intel Core i7-12700 (20 threads) for NR with BLAS set to one thread per
worker to avoid oversubscription.

\textbf{Batched:} Yes---streaming micro-batches of 4,096 cases filling GPU
memory without overflow.

\textbf{Why this is the most honest benchmark in the literature:}
This is the only paper that uses a genuinely competitive NR baseline
(multi-worker, vectorised, analytically assembled Jacobian) rather than
single-threaded Python NR or pandapower.  The resulting $2$--$5\times$
speedup with a modest GPU is far more realistic than the headline numbers
from other papers.  Critically, the paper explicitly states: ``In
single-scenario runs, NR is competitive for $N < 200$''---directly explaining
why our 95-node result shows the GNN behind.  Figure~\ref{fig:pignn_timing}
makes this crossover visually explicit: the left panel shows NR beating all
PIGNN variants below roughly 200 buses, after which PIGNN's GPU parallelism
decisively takes over.

\subsection{B\"{o}ttcher et al.\ (2023)}

\textbf{Citation:} Luis B\"{o}ttcher et al.
``Solving AC Power Flow with Graph Neural Networks under Realistic
Constraints.''
In \textit{2023 IEEE Belgrade PowerTech}, pp.~1--7.
DOI:~10.1109/PowerTech55446.2023.10202246.
arXiv: \url{https://arxiv.org/abs/2204.07000}

\textbf{Networks tested:} Synthetic MV distribution grids of 20--500 buses;
also SimBench MV (100-bus) and a real MV grid (120-bus).

\textbf{Speed claims (Figure~4, Section~V):}
GNN inference time is nearly constant across grid sizes (batch of 50 grids);
Newton--Raphson (PYPOWER) scales up with problem size.
The crossover point where GNN becomes faster is approximately 200--300 buses
(visual reading of Figure~4).

\textbf{Hardware:} NVIDIA Tesla V100 16\,GB GPU for GNN; PYPOWER on CPU.

\textbf{Batched:} Yes---50 grids batched per evaluation; additionally the
method uses 50 recurrent iterations $\times$ 10 random initialisations
(500 forward passes per problem), all included in reported time.

\textbf{Critical caveat:}  The 500 forward passes per inference scenario
inflate GNN wall time substantially---this is an unsupervised approach that
trades many forward passes for physics consistency without NR ground-truth
training.  At 95 nodes, NR would still be faster based on Figure~4's trend.

\subsection{Hybrid GNN-IZR (Shamseldein, 2025)}

\textbf{Citation:} Mohamed Shamseldein.
``A Hybrid GNN-IZR Framework for Fast and Empirically Robust AC Power Flow
Analysis in Radial Distribution Systems.''
arXiv preprint, arXiv:2510.04264, October 2025.
\url{https://arxiv.org/abs/2510.04264}

\textbf{Networks tested:} IEEE 33-bus radial distribution system only
(7,500 stressed test scenarios).

\textbf{Speed result (Figure~6(a); Table~III):}
The hybrid GNN-IZR framework has a mean solution time of ${\approx}3.794\,\text{ms}$,
while standard NR (PYPOWER) averages ${\approx}1.382\,\text{ms}$---making the
hybrid \textcolor{warnred}{\textbf{${\approx}2.7\times$ slower}} than NR.
The raw GNN-only path has a median of ${\approx}3.01\,\text{ms}$.

\textbf{Hardware:} GNN on GPU; IZR solver Numba JIT-compiled on CPU;
NR baseline from PYPOWER on CPU.

\textbf{Batched:} No---single-scenario per-case inference.

\textbf{Why this is the most relevant paper to our work:}
This is the only paper in this group where the GNN is demonstrably slower
than the traditional solver on a small distribution feeder (33-bus), closely
mirroring our finding of ${\approx}11.7\times$ slower on CPU.  The paper implicitly
acknowledges that for small radial distribution networks, GNN inference
overhead exceeds iterative solver time, and pivots the value proposition
from speed to robustness (0\% convergence failures vs.\ 13.11\% for the
pure GNN approach on stressed scenarios).  \textbf{This is exactly the
situation in this work}: our 95-node feeder is small enough that OpenDSS
wins on raw speed, but the GNN surrogate offers other potential benefits
(differentiability, topology generalisation, GPU throughput at batch scale).

\subsection{Synthesis: Why the Present Results Are Consistent}

Three structural factors explain every entry in Table~\ref{tab:lit} and
position the present benchmark precisely within the literature.

\begin{description}[style=unboxed, leftmargin=0pt, labelwidth=0pt]

  \item[\textbf{Network size is the dominant variable.}]
  PowerFlowMultiNet shows speedups scaling from $15\times$ (13-bus) to
  $38\times$ (906-bus).  PIGNN-Attn-LS explicitly finds NR competitive below
  ${\approx}200$ buses.  B\"{o}ttcher et al.'s crossover is at 200--300 buses.
  Our \textbf{95-node feeder sits squarely in the regime where traditional
  solvers are expected to win.}

  \item[\textbf{GPU-vs-CPU comparisons inflate speedup ratios.}]
  Every paper reporting large speedups ($15$--$38\times$, $8{,}400\times$)
  compares GPU inference against single-threaded CPU solvers.  The only paper
  using a strong multi-worker CPU baseline (PIGNN-Attn-LS, 20 NR workers)
  reports a much more modest $2$--$5\times$ speedup, and only at scales above
  200 buses.  Our finding of ${\approx}11.7\times$ slower on CPU is consistent with this:
  CPU GNN inference is inherently disadvantaged by the lack of parallelism
  that makes GPUs attractive.

  \item[\textbf{Batching is essential for GNN competitiveness.}]
  Our finding that GPU + batching narrows the gap to $1.6\times$ slower aligns
  with the literature.  The GNN-IZR paper, which uses single-scenario inference
  on a 33-bus system, shows the GNN at ${\approx}3\,\text{ms}$ vs.\ NR at
  ${\approx}1.4\,\text{ms}$---the GNN is slower.  Without batching, GNN
  per-scenario latency exceeds solver time for small networks.

\end{description}

% =============================================================================
\section{Roadmap for GNN-Beats-OpenDSS at Batch Size One}
\label{sec:roadmap}
% =============================================================================

A GNN surrogate can plausibly outperform OpenDSS in single-scenario
sequential mode, but only on distribution feeders with roughly
2,000--5,000+ buses.  The core challenge is that OpenDSS uses an
extraordinarily fast fixed-point current-injection solver with KLU sparse
factorisation that scales nearly linearly for radial networks, while a
batch-size-1 GPU GNN forward pass carries ${\approx}1$--$3$\,ms of fixed
overhead from kernel launches and CPU--GPU dispatch.  The crossover occurs
when network size pushes OpenDSS solve time above that fixed overhead.
The bottleneck to beat is therefore not OpenDSS's algorithm---which is
near-optimal for radial networks---but the PyTorch/CUDA dispatch overhead,
making inference optimisation the single highest-leverage investment.

% ── OpenDSS scaling ─────────────────────────────────────────────────────────
\subsection{OpenDSS Is Deceptively Fast on Radial Networks}

OpenDSS does not use Newton--Raphson.  Its default mode is a
\textbf{fixed-point current-injection method} that builds the system
admittance matrix once, factors it via the KLU sparse direct solver, and
then iterates by updating current injections and performing
forward/backward substitution.  For sequential time-series (QSTS)
simulations, convergence typically requires only \textbf{2 iterations per
timestep} because the previous solution provides an excellent warm start.
The KLU factorisation for radial (tree) networks produces essentially zero
fill-in, yielding $O(n)$ factorisation and $O(n)$ solve per iteration.

Calibrating against the 95-node feeder benchmarked in this work
(${\approx}0.2$\,ms per solve), the estimated per-solve times for larger
feeders are:

\begin{table}[ht]
\centering
\renewcommand{\arraystretch}{1.25}
\begin{tabular}{
  >{\raggedright\arraybackslash}p{3.2cm}
  >{\centering\arraybackslash}p{3.6cm}
  >{\centering\arraybackslash}p{3.6cm}
}
\toprule
\rowcolor{headerblue}
\textcolor{white}{\textbf{Network size}}
  & \textcolor{white}{\textbf{OpenDSS (no controls)}}
  & \textcolor{white}{\textbf{OpenDSS (with controls)}} \\
\midrule
\rowcolor{rowblue}
89 buses (this work)      & ${\approx}0.2$\,ms & ${\approx}0.3$--$0.5$\,ms \\
500 buses                  & ${\approx}0.5$--$1.5$\,ms & ${\approx}1$--$4$\,ms \\
\rowcolor{rowblue}
1,000 buses                & ${\approx}1$--$3$\,ms & ${\approx}3$--$8$\,ms \\
3,000 buses                & ${\approx}3$--$8$\,ms & ${\approx}8$--$25$\,ms \\
\rowcolor{rowblue}
5,000 buses                & ${\approx}5$--$15$\,ms & ${\approx}15$--$50$\,ms \\
8,500 nodes (IEEE 8500)   & ${\approx}10$--$30$\,ms & ${\approx}30$--$100$\,ms \\
\bottomrule
\end{tabular}
\caption{Estimated OpenDSS per-solve times calibrated from the 95-node
  benchmark.  Voltage regulators, capacitor banks, and smart-inverter
  controls add outer control iterations that multiply solve time by
  $2$--$5\times$.}
\label{tab:opendss_scaling}
\end{table}

The critical observation: PowerFlowNet's headline ``$48$--$145\times$
speedup'' is measured against pandapower Newton--Raphson, which is
${\approx}117\times$ slower than OpenDSS.  Against OpenDSS specifically,
the advantage shrinks dramatically.

% ── Fixed overhead floor ─────────────────────────────────────────────────────
\subsection{The Fixed Overhead Floor at Batch Size One}

At batch size~1 on a small-to-medium graph, GNN inference is dominated by
fixed overhead rather than computation.  A typical PyTorch Geometric
forward pass with 2--3 message-passing layers launches ${\approx}10$--$30$
CUDA kernels.  Python/ATen operator dispatch, memory management, and
CPU--GPU synchronisation push the effective per-launch cost to
${\approx}20$--$200\,\mu$s, giving a total fixed overhead in
\textbf{eager mode of ${\approx}1.5$--$3.0$\,ms} regardless of graph
size.

The present benchmark confirms this directly.  At 95 nodes, the GNN on GPU
with batching runs at ${\approx}1.73$\,ms per step; at batch size~1 it
runs at ${\approx}2.38$\,ms---a difference of only $0.65$\,ms despite
$288\times$ fewer nodes being processed.  The computation is negligible;
the overhead is not.  The good news: this overhead is compressible.

% ── Three optimizations ──────────────────────────────────────────────────────
\subsection{Three Optimisation Techniques That Change the Crossover Math}

\textbf{\texttt{torch.compile} with \texttt{mode=\char`\"reduce-overhead\char`\"}:}
Fuses multiple PyG message-passing operations into fewer CUDA kernels,
applies constant folding, and automatically employs CUDA Graphs.
PyG~2.5+ explicitly supports this.  Documented speedups on
GCN/GraphSAGE/GIN architectures reach \textbf{up to $3\times$}, reducing
overhead from ${\approx}2$\,ms to ${\approx}0.7$--$1.0$\,ms.

\textbf{CUDA Graphs:}
Record the entire kernel sequence on the first forward pass and replay it
as a single GPU operation on subsequent calls, eliminating per-kernel CPU
dispatch.  This requires fixed input tensor shapes---perfectly satisfied
by a fixed-topology distribution network where only node features ($P$,
$Q$ injections) change between timesteps.  Combined with
\texttt{torch.compile}, overhead can drop below \textbf{0.3--0.5\,ms}.

\textbf{FP16 mixed precision:}
\texttt{torch.autocast} engages Tensor Cores on V100/A100/RTX-series GPUs,
roughly doubling throughput for matrix multiplications within
message-passing layers.  The benefit is modest at small graph sizes where
compute is not the bottleneck, but increasingly significant above
${\approx}2{,}000$ nodes.

With all three optimisations applied, the estimated GNN batch=1 inference
time compresses to approximately \textbf{0.3--0.5\,ms at 95 nodes} and
\textbf{0.5--2.0\,ms at 5,000+ nodes}.  This shifts the crossover
substantially.

% ── Crossover analysis ───────────────────────────────────────────────────────
\subsection{Calibrated Crossover Analysis}

Using the 95-node data point as calibration anchor (OpenDSS
${\approx}0.2$\,ms, GNN batch=1 GPU ${\approx}1.5$--$3$\,ms in eager
mode), the crossover bus count depends on both OpenDSS control complexity
and GNN optimisation level:

\begin{table}[ht]
\centering
\renewcommand{\arraystretch}{1.25}
\begin{tabular}{
  >{\raggedright\arraybackslash}p{4.6cm}
  >{\centering\arraybackslash}p{2.2cm}
  >{\centering\arraybackslash}p{2.0cm}
  >{\centering\arraybackslash}p{2.8cm}
}
\toprule
\rowcolor{headerblue}
\textcolor{white}{\textbf{Scenario}}
  & \textcolor{white}{\textbf{OpenDSS scaling}}
  & \textcolor{white}{\textbf{GNN overhead}}
  & \textcolor{white}{\textbf{Crossover point}} \\
\midrule
\rowcolor{rowblue}
Simple radial, unoptimised GNN & $O(n^{1.0})$ & ${\approx}2$\,ms & ${\approx}900$ buses \\
Simple radial, optimised GNN   & $O(n^{1.0})$ & ${\approx}0.5$\,ms & ${\approx}220$ buses \\
\rowcolor{rowblue}
Moderate controls, unoptimised & $O(n^{1.2})$ & ${\approx}2$\,ms & 500--1,500 buses \\
Moderate controls, optimised   & $O(n^{1.2})$ & ${\approx}0.5$\,ms & 200--500 buses \\
\rowcolor{rowblue}
Complex controls, unoptimised  & $O(n^{1.5})$ & ${\approx}2$\,ms & 300--800 buses \\
\bottomrule
\end{tabular}
\caption{Estimated crossover bus count at which a batch-size-1 GNN becomes
  faster than OpenDSS, as a function of feeder control complexity and GNN
  optimisation level.}
\label{tab:crossover}
\end{table}

\noindent\textbf{The sweet spot for a convincing demonstration is a feeder
in the 2,000--5,000 bus range.}  At this size, even an unoptimised GNN is
within striking distance, and an optimised one should be $2$--$5\times$
faster than OpenDSS.  At 8,500 nodes, the advantage grows to an estimated
$5$--$20\times$ depending on control complexity.

A critical tension exists: a topologically simple feeder (few controls)
minimises GNN input features but also minimises OpenDSS solve time.
The optimal strategy is a feeder that is \textbf{large} (to push OpenDSS
above the overhead threshold) but \textbf{topologically simple} (to keep
the GNN minimal), while accepting that some voltage-dependent controls
actually help because they slow OpenDSS disproportionately through outer
control iterations.

% ── Test systems ─────────────────────────────────────────────────────────────
\subsection{Available Test Systems Ranked by Suitability}

\textbf{Tier 1 (best candidates).}
The \textbf{IEEE 8500-node test feeder} (4,876 buses / 8,531 nodes) is
the strongest candidate.  It is included in every OpenDSS installation
(\texttt{IEEETestCases/8500-Node/}), is the most widely benchmarked large
distribution system, and has 4 voltage regulators, 4 capacitor banks, and
${\approx}1{,}100$ centre-tapped transformers---which significantly slow
OpenDSS through outer control iterations.  A balanced single-phase
equivalent preserves the large topology while reducing GNN input
dimensionality.

\textbf{NREL Smart-DS feeders} are synthetic city-scale distribution
feeders (San Francisco, Greensboro, Austin) in native OpenDSS format,
available at \url{https://data.openei.org/submissions/2981}.  Individual
feeders range from hundreds to tens of thousands of nodes and can be
selected at any desired size.

\textbf{Tier 2 (good candidates).}
\textbf{EPRI J1} at 3,434 buses / 4,245 nodes is the largest EPRI test
circuit included with OpenDSS and already models distributed generation.
\textbf{EPRI Ckt7} has 5,694 customers in a compact, heavily loaded
topology.  Both are unbalanced and in native OpenDSS format.

\textbf{Tier 3 (supplementary).}
The \textbf{IEEE 906-bus European LV feeder} is purely radial with no
voltage regulators or capacitors, making it topologically simple but
possibly too small (${\approx}906$ buses) for a convincing crossover with
an unoptimised GNN.
\textbf{PNNL taxonomy feeders} (up to ${\approx}850$--$1{,}630$ nodes in
OpenDSS format) are useful for intermediate-size validation but below the
primary demonstration threshold.

% ── Minimal GNN architecture ─────────────────────────────────────────────────
\subsection{Minimal GNN Architecture for the Target Feeder}

For a large simplified distribution feeder, the GNN can be remarkably
lightweight.  The minimum viable architecture is:

\begin{itemize}
  \item \textbf{Input features per node:} 2 ($P$ and $Q$ injections).
        Line impedances are fixed for a given topology and can be
        pre-encoded as fixed edge attributes.  Output: 2 features
        (voltage magnitude and angle).  For an unbalanced per-phase model:
        6 inputs / 6 outputs per bus.

  \item \textbf{Architecture:} 2--3 message-passing layers, hidden
        dimension $h = 32$--$64$.  TAGConv with $K=3$ gives an effective
        9-hop receptive field in 3 layers---sufficient because voltage
        drops in distribution systems are primarily local phenomena, as
        confirmed by PowerFlowNet for networks up to 6,470 buses.
        Total parameters: ${\approx}10{,}000$--$50{,}000$, compared to
        PowerFlowNet-Medium's 357,000.

  \item \textbf{Edge features:} 2 ($r$, $x$ per line), pre-loaded as
        persistent GPU tensors---zero runtime overhead.

  \item \textbf{Estimated inference time:}  ${\approx}2$--$4$\,ms in
        eager mode for a 3,000--8,500 node graph at batch=1; compresses
        to \textbf{${\approx}0.5$--$1.5$\,ms with
        \texttt{torch.compile} + CUDA Graphs}.
\end{itemize}

% ── Concrete roadmap ─────────────────────────────────────────────────────────
\subsection{Revised Roadmap: Single-Scenario Sequential Inference on the Current Feeder}
\label{subsec:revised_roadmap}

The roadmap below is designed around the actual constraint: the same
95-node IEEE~34-bus feeder, a battery controller that enforces sequential
timesteps, and a target of beating OpenDSS wall-clock time at batch
size~1.  Steps are ordered to win the easiest gains first and to minimise
interaction with distribution-system complexities (voltage-dependent loads,
regulators, capacitor banks) until the inference pipeline is already fast.

\medskip
\noindent\textbf{Phase 1 --- Vectorise the feature-building pipeline (1 week).}

The first bottleneck that can be eliminated without touching the model is
the Python loop inside \texttt{build\_gnn\_x}.  Currently this function
runs per-timestep in pure Python, costing ${\approx}0.34$\,ms/step on GPU
non-batched.  The fix is to pre-stack all time-varying inputs as a
persistent numpy or tensor array at the start of the simulation and index
into it with a single slice per step.  Separately, because the feeder
topology is fixed, \texttt{edge\_index} and all edge features
(impedances, line lengths, phase configuration) should be converted to
persistent GPU tensors at startup and never re-transferred.  Similarly,
\texttt{tensor\_data\_creation} should be replaced by a direct tensor
slice and offset-addition rather than \texttt{Batch.from\_data\_list()},
eliminating the sequential Python concatenation loop that currently costs
${\approx}0.15$\,ms/step.  After this phase the pipeline overhead
stages should be negligible ($<0.05$\,ms/step combined), leaving the
model forward pass as the only remaining cost.

\medskip
\noindent\textbf{Phase 2 --- Establish an MLP baseline (1--2 weeks).}

Before investing further in GNN optimisation, train a shallow MLP on the
same dataset and inputs.  The MLP takes the full node-feature matrix
flattened into a single vector (all 95 buses concatenated) and outputs
predicted voltages for all buses.  Conrad et al.~\cite{Conrad2026}
show that on CPU this architecture is roughly two orders of magnitude
faster than a GNN at the sample sizes relevant here ($N \leq 288$).
Because the feeder topology is fixed, the MLP's lack of topological
inductive bias is not a disadvantage---it never needs to generalise to
unseen topologies.  Moreover, because distribution feeder buses are
heterogeneous (different load types, regulator buses, capacitor buses),
the MLP's fully independent per-input weights may actually learn
bus-specific behaviour better than the GNN's shared message-passing
rule.  The MLP baseline therefore provides both a speed ceiling and an
accuracy sanity check: if the MLP already beats OpenDSS on GPU and
matches GNN accuracy, the GNN architecture is unnecessary for this
feeder.

\medskip
\noindent\textbf{Phase 3 --- Optimise the GPU inference pipeline (1 week).}

Apply the following in sequence, benchmarking after each step:

\begin{enumerate}[itemsep=2pt]
  \item \textbf{\texttt{torch.inference\_mode()}} instead of
        \texttt{no\_grad}---saves autograd bookkeeping overhead.
  \item \textbf{FP16 autocast}---halves memory bandwidth for weight reads;
        safe for inference if trained in FP32.
  \item \textbf{Pre-pinned input buffer}---allocate the node-feature tensor
        with \texttt{pin\_memory=True} so CUDA can DMA directly from
        page-locked RAM.  Transfer with
        \texttt{.to(device, non\_blocking=True)}.
  \item \textbf{Dual CUDA streams}---one stream runs the forward pass for
        step $t$ while a second stream copies features for step $t{+}1$,
        overlapping transfer and compute entirely.
  \item \textbf{\texttt{torch.compile(model, mode=\char`\"reduce-overhead\char`\",
        dynamic=False)}}---fuses kernels, applies constant folding, and
        wraps the repeated identical call-graph into a CUDA Graph,
        eliminating the ${\approx}10\,\mu$s per-kernel Python dispatch
        overhead.  Expected reduction: forward pass from
        ${\approx}1.9$\,ms/step to ${\approx}0.3$--$0.5$\,ms/step.
\end{enumerate}

\noindent Target after Phase~3: end-to-end inference below $0.5$\,ms/step,
which would place the GNN (or MLP) faster than OpenDSS wall-clock
($1.38$\,ms/step) with margin to spare.

\medskip
\noindent\textbf{Phase 4 --- Handle distribution-system inputs correctly (1--2 weeks).}

This phase addresses the distribution-specific complexities that can
silently degrade accuracy if handled incorrectly.

\emph{Voltage-dependent (ZIP) loads.}  Never use post-solve $P$ and $Q$
values as model inputs---these are only known \emph{after} the power flow
converges and create a circular dependency (the PowerFlowMultiNet
problem, Section~\ref{sec:multinet}).  Use nominal $P$ and $Q$ values
(pre-solve, as specified in the load profile) and the ZIP coefficients as
static node features.  The model then learns the direct mapping from
nominal load demand to converged bus voltages; the voltage-correction
effect is implicitly absorbed into the training targets from OpenDSS.
Validate that inference error does not increase on timesteps with large
voltage deviations from nominal (where ZIP correction is most significant).

\emph{Voltage regulators.}  The regulator tap position after convergence
is a discrete state that affects downstream voltages strongly.  Add the
per-tap-step voltage ratio of each regulator as a node or edge feature,
sampled from the OpenDSS state \emph{at the start of each timestep}
(before the solve).  This is available without running the solver.
Because ControlIter $\approx 2.33$ in this feeder, tap positions do
change within a timestep; the model must learn to predict the post-control
voltage from the pre-control tap state.  Validate on timesteps where
tap changes occur.

\emph{Capacitor banks.}  Add the switched state (in/out) of each
capacitor bank as a binary node feature.  Same principle: use the
pre-solve state, validate on switching timesteps.

\emph{Three-phase imbalance.}  The IEEE~34-bus feeder is unbalanced.
If the model operates on a per-phase graph (as the current GNN3 Best7
does), ensure phase identity is encoded in node features.  If operating
on a single-phase equivalent, validate that phase-averaging does not
introduce systematic error on heavily unbalanced segments.

\medskip
\noindent\textbf{Phase 5 --- End-to-end evaluation and decision point (1 week).}

Run the full battery controller simulation loop (288 timesteps, 24-hour
profile) with OpenDSS replaced by the optimised model.  Measure
wall-clock time per timestep including feature construction, host$\to$device
transfer, model forward, and result extraction.  Compare against
the current OpenDSS baseline of $1.38$\,ms/step.

At this point one of three outcomes is expected:

\begin{description}[itemsep=3pt]
  \item[\textbf{Success (model faster, accuracy acceptable).}]
    Proceed to deployment.  The MLP variant is the likely winner here
    given Conrad et al.'s evidence.

  \item[\textbf{Close but not faster.}]
    The remaining gap is most likely the model forward pass.  The next
    lever is architecture reduction: halve the hidden dimension (128
    $\to$ 64 or 32) and retrain.  Each halving reduces the forward
    FLOP count by $4\times$.  Accept the accuracy penalty only if
    it remains within operational tolerance.

  \item[\textbf{Still substantially slower.}]
    The 95-node feeder is below the single-scenario crossover point
    (${\approx}200$ buses) identified in the literature.  Pivot the
    use case: abandon real-time single-scenario replacement and instead
    use the model for high-throughput \emph{scenario screening}
    (thousands of independent load profiles evaluated in parallel),
    where the GPU advantage is decisive and the sequential dependency
    of the battery controller does not apply.
\end{description}

% =============================================================================
\section{Practical Implementation Checklist: IEEE 8500-Node Feeder Pipeline}
\label{sec:checklist}
% =============================================================================

This section is a sequential, hands-on checklist for building the full
GNN-beats-OpenDSS pipeline from scratch on the IEEE~8500-node feeder.
Follow the steps in order; each step has a clear pass/fail criterion before
you continue to the next.

% ---------------------------------------------------------------------------
\subsection{Step 1: Locate or Install the IEEE 8500-Node Feeder}
\label{subsec:step1}
% ---------------------------------------------------------------------------

\noindent\textbf{What to do.}
The IEEE~8500-node feeder ships as a set of \texttt{.dss} files
(entry point: \texttt{Master.dss}) and is distributed with the official
OpenDSS test cases.  Run the following search in any Python environment that
has \texttt{pathlib} available to find it on the machine:

\begin{verbatim}
import pathlib
roots = [
    pathlib.Path(r"C:\Users\alita\OneDrive\Desktop\quest-ssim"),
    pathlib.Path(r"C:\Users\alita\Documents"),
    pathlib.Path(r"C:\Users\alita\Desktop"),
    pathlib.Path(r"C:\Users\alita\Downloads"),
    pathlib.Path(r"C:\OpenDSS"),
    pathlib.Path(r"C:\Program Files\OpenDSS"),
]
for root in roots:
    if root.exists():
        for f in root.rglob("*.dss"):
            if "8500" in str(f).lower():
                print(f)
\end{verbatim}

\noindent\textbf{If not found --- Option A (recommended).}
Download the feeder directly from the official OpenDSS GitHub mirror:

\begin{verbatim}
# In a terminal / Anaconda Prompt (no admin rights needed):
pip install dss-python          # ships the full IEEE test suite
python -c "import dss; print(dss.__file__)"
# Navigate to that folder -> IEEETestCases -> 8500-Node
\end{verbatim}

\noindent\textbf{If not found --- Option B.}
Clone or download the zip from the OpenDSS SourceForge repository:
\url{https://sourceforge.net/p/electricdss/code/HEAD/tree/trunk/Distrib/IEEETestCases/8500-Node/}
and extract to a known path (e.g.\ \texttt{C:\textbackslash OpenDSS\textbackslash 8500-Node\textbackslash}).

\noindent\textbf{Pass criterion.}
You have a local path to \texttt{Master.dss} that can be stored in a
variable \texttt{DSS\_8500} and passed to \texttt{dss.run\_command('Redirect ...')}.

% ---------------------------------------------------------------------------
\subsection{Step 2: Profile the OpenDSS Solve Time}
\label{subsec:step2}
% ---------------------------------------------------------------------------

\noindent\textbf{What to do.}
In the \textit{try for 8500 node} Jupyter notebook, set \texttt{DSS\_8500}
to the path found in Step~1 and run the profiling cell (already written in
the notebook).  The cell loads the feeder, runs 50 snapshot solves, and
reports median/mean/min/max solve time in milliseconds.

\noindent\textbf{What to check.}
\begin{itemize}
  \item \texttt{Converged: True} --- if False, the feeder has a load or
    voltage issue; check that \texttt{Redirect} picked up all dependency
    files correctly.
  \item Buses: ${\approx}4{,}876$, node-phases: ${\approx}8{,}531$ (hence
    ``8500-node'').
  \item Median solve time should be in the range of \textbf{5--30\,ms} per
    snapshot on a typical laptop.  This is the GNN target to beat.
  \item Record \texttt{ControlIterations} and \texttt{Iterations}
    (power-flow iterations).  A high control-iteration count
    ($>$5) means the feeder has many regulators or capacitors switching
    and will be harder to surpass; a low count ($<$3) is favourable.
\end{itemize}

\noindent\textbf{Pass criterion.}
Median solve time $\geq 2$\,ms per snapshot.  If it is below 2\,ms the
feeder may still be too fast; if above 5\,ms the GNN will have a very
comfortable margin.

% ---------------------------------------------------------------------------
\subsection{Step 3: Generate the Training Dataset}
\label{subsec:step3}
% ---------------------------------------------------------------------------

\noindent\textbf{What to do.}
Write a standalone script \texttt{generate\_dataset\_8500.py} that:

\begin{enumerate}[itemsep=2pt]
  \item Loads the feeder with \texttt{opendssdirect} and sets
    \texttt{dss.Solution.Mode(1)} (snapshot).
  \item Generates $N=5{,}000$--$10{,}000$ random load scenarios by scaling
    each load's $P$ and $Q$ independently, drawn from a truncated normal
    distribution centred on the nominal value with $\sigma = 20\%$.
    Use \texttt{dss.Loads.kW()} and \texttt{dss.Loads.kvar()} setters.
  \item Records the \emph{pre-solve} state vector $\mathbf{x}$ for each
    scenario: nominal $P$, $Q$ for every load bus; regulator tap positions
    (read with \texttt{dss.RegControls.TapNumber()}); capacitor switch
    states (\texttt{dss.Capacitors.States()}).
  \item Solves and records the \emph{post-solve} voltage magnitudes and
    angles for all node-phases as the target vector $\mathbf{y}$.
  \item Saves the full dataset as compressed numpy arrays:
    \texttt{X\_8500.npy} (shape: $N \times F$) and
    \texttt{Y\_8500.npy} (shape: $N \times 2 \times n_{\text{nodes}}$
    for magnitude and angle).
\end{enumerate}

\noindent\textbf{Key implementation note.}
Do \emph{not} use post-solve $P$ and $Q$ as model inputs.  Use only the
pre-solve nominal demand values --- these are causally prior to the solution
and avoid the circular-dependency problem described in
Section~\ref{sec:multinet}.

\noindent\textbf{Expected runtime.}
At $\approx$10\,ms per solve, 10\,000 scenarios takes about 100\,s.  Run
in a single Python process with no parallelism needed.

\noindent\textbf{Pass criterion.}
\texttt{X\_8500.npy} and \texttt{Y\_8500.npy} saved and loadable.  Spot-check
5 random rows: all voltage magnitudes should be in the range 0.85--1.10\,pu
and all scenarios should have converged (check by saving a convergence flag
alongside each row and asserting all are \texttt{True}).

% ---------------------------------------------------------------------------
\subsection{Step 4: Build the Graph Structure}
\label{subsec:step4}
% ---------------------------------------------------------------------------

\noindent\textbf{What to do.}
Write \texttt{build\_graph\_8500.py} that constructs the static PyG
\texttt{Data} object for the feeder (this is computed once and cached):

\begin{enumerate}[itemsep=2pt]
  \item Load the feeder (without solving) and use
    \texttt{dss.Lines.AllNames()}, \texttt{dss.Lines.Bus1()},
    \texttt{dss.Lines.Bus2()}, \texttt{dss.Lines.R1()},
    \texttt{dss.Lines.X1()}, \texttt{dss.Lines.Length()} to build
    \texttt{edge\_index} (shape $2 \times E$) and \texttt{edge\_attr}
    (impedance + length per edge).
  \item Assign each bus a node ID (integer) matching the row index in
    \texttt{X\_8500.npy}.  Store the mapping as a dict.
  \item Save \texttt{edge\_index.pt}, \texttt{edge\_attr.pt}, and
    the bus-name-to-index mapping as a JSON file.
  \item Load these to GPU at startup and keep them there for the full
    simulation:
    \begin{verbatim}
edge_index = torch.load("edge_index.pt").to("cuda")
edge_attr  = torch.load("edge_attr.pt").to("cuda")
    \end{verbatim}
\end{enumerate}

\noindent\textbf{Pass criterion.}
\texttt{edge\_index} has shape $(2, E)$ with $E \approx 8{,}000$--$9{,}000$
(roughly one edge per node in a radial feeder).  Visualise the degree
distribution: most buses should have degree 1 or 2 (radial), with a handful
of junction buses at degree 3--4.

% ---------------------------------------------------------------------------
\subsection{Step 5: Train the MLP Baseline}
\label{subsec:step5}
% ---------------------------------------------------------------------------

\noindent\textbf{What to do.}
Before building the GNN, train an MLP as a speed and accuracy baseline.
Write \texttt{train\_mlp\_8500.py}:

\begin{enumerate}[itemsep=2pt]
  \item Load \texttt{X\_8500.npy} and \texttt{Y\_8500.npy}.  Split 80/10/10
    train/val/test.
  \item Normalise: z-score each feature column.  Save the mean and std
    vectors --- they will be applied at inference time.
  \item Architecture: \texttt{Linear($F$, 512) $\to$ ReLU $\to$
    Linear(512, 512) $\to$ ReLU $\to$ Linear(512, $2 \times n_{\text{nodes}}$)}.
    Two hidden layers of width 512 is a safe starting point; the feeder has
    ${\approx}8{,}531$ node-phases so the output layer is large.
  \item Loss: MSE on voltage magnitude; optionally add MSE on angle with a
    lower weight (0.1) since angles have much smaller variance.
  \item Train for 100 epochs with Adam, lr $= 10^{-3}$, reduce on plateau.
  \item Evaluate on test set: report mean absolute error (MAE) in volts
    (multiply pu error by nominal voltage).
\end{enumerate}

\noindent\textbf{Pass criterion.}
Test MAE $< 0.005$\,pu on voltage magnitude ($< 0.6$\,V at 120\,V base).
This is the accuracy bar --- the GNN must match or beat it.

% ---------------------------------------------------------------------------
\subsection{Step 6: Train the GNN}
\label{subsec:step6}
% ---------------------------------------------------------------------------

\noindent\textbf{What to do.}
Write \texttt{train\_gnn\_8500.py} using the same dataset.  Recommended
architecture for a first pass:

\begin{enumerate}[itemsep=2pt]
  \item \textbf{Node embedding}: linear projection of the node feature vector
    to hidden dimension $H = 64$.  Use node-type one-hot encoding (load bus,
    regulator bus, capacitor bus, source bus) as additional static features
    concatenated to the per-timestep features.
  \item \textbf{Message passing}: 3--4 layers of
    \texttt{GCNConv} or \texttt{SAGEConv} with $H=64$, ReLU activation,
    residual connections after each layer.
  \item \textbf{Readout}: per-node \texttt{Linear(64, 2)} (magnitude + angle).
  \item Same loss, optimiser, and train/val/test split as the MLP.
\end{enumerate}

\noindent\textbf{Why $H=64$?}  The 8500-node feeder is large enough that
GNN message passing provides genuine value (information propagates over
${\approx}30$--50 hops in a radial feeder of this size), but $H=128$ would
increase the forward-pass FLOP count $4\times$ for marginal accuracy gain.
Start at 64 and increase only if the accuracy target is not met.

\noindent\textbf{Pass criterion.}
Test MAE $\leq$ MLP baseline.  If the GNN is more than 2$\times$ worse than
the MLP, the MLP is the better model for this feeder.

% ---------------------------------------------------------------------------
\subsection{Step 7: Benchmark Inference Speed}
\label{subsec:step7}
% ---------------------------------------------------------------------------

\noindent\textbf{What to do.}
Write \texttt{benchmark\_inference\_8500.py}.  For both the MLP and GNN:

\begin{enumerate}[itemsep=2pt]
  \item Load model weights and move to GPU.  Call
    \texttt{torch.compile(model, mode="reduce-overhead", dynamic=False)}
    and run 20 warm-up steps (the first call triggers compilation).
  \item Pre-build the full feature matrix \texttt{X} (all $N$ rows) and keep
    on CPU with \texttt{pin\_memory=True}.
  \item Simulate the sequential battery-controller loop:
    for each timestep $t = 1 \ldots N$, copy one row of \texttt{X} to GPU
    with \texttt{non\_blocking=True}, run the model forward, move the voltage
    result back to CPU.  Time each step with \texttt{time.perf\_counter()}.
  \item Report median, mean, p95 wall-clock time per step.
\end{enumerate}

\noindent\textbf{What to compare.}
\begin{itemize}
  \item OpenDSS median solve time (from Step~2).
  \item MLP median inference time per step.
  \item GNN median inference time per step.
\end{itemize}

\noindent\textbf{Pass criterion.}
At least one of (MLP, GNN) has median inference time \textbf{below} the
OpenDSS median solve time from Step~2.  Given that the 8500-node feeder
is well above the ${\approx}200$-bus crossover identified in the PIGNN
paper, this outcome is strongly expected.

% ---------------------------------------------------------------------------
\subsection{Step 8: End-to-End Validation in the Battery Controller}
\label{subsec:step8}
% ---------------------------------------------------------------------------

\noindent\textbf{What to do.}
Replace the \texttt{dss.Solution.Solve()} call inside the existing battery
controller loop with the trained model.  Run the full 24-hour simulation
(288 timesteps) and compare:

\begin{enumerate}[itemsep=2pt]
  \item \textbf{Wall-clock time}: total simulation time with OpenDSS vs.\
    with the model.
  \item \textbf{Voltage accuracy}: plot predicted vs.\ true voltage at
    a representative set of buses across all 288 timesteps.
  \item \textbf{Battery dispatch accuracy}: compare the battery charge/discharge
    schedule produced by the controller under predicted voltages vs.\ true
    voltages.  Assess whether voltage prediction errors propagate into
    materially different control decisions.
\end{enumerate}

\noindent\textbf{Pass criterion.}
Total simulation time with model $<$ total simulation time with OpenDSS,
\emph{and} battery dispatch decisions differ by $<5\%$ from the OpenDSS
ground truth.

% =============================================================================
% References
% =============================================================================
\begin{thebibliography}{99}

\bibitem{lightsim2grid}
B.~Donnot,
``lightsim2grid: A fast simulator for the Grid2op platform,''
\textit{lightsim2grid documentation, v0.9.2},
RTE France, 2023.
\url{https://lightsim2grid.readthedocs.io/en/v0.9.2/benchmarks_grid_sizes.html}

\bibitem{Davis2010}
T.~A.~Davis and E.~Palamadai Natarajan,
``Algorithm 907: KLU, a direct sparse solver for circuit simulation
problems,''
\textit{ACM Trans.\ Math.\ Softw.}, vol.~37, no.~3, pp.~36:1--36:17, 2010.
DOI:~10.1145/1824801.1824814.

\bibitem{pytorch_overhead}
PyTorch contributors,
``Large overhead (7 microseconds) for PyTorch operation,''
\textit{PyTorch GitHub Issue \#41383}, 2020.
\url{https://github.com/pytorch/pytorch/issues/41383}

\bibitem{pyg_compile}
PyTorch Geometric contributors,
``Compiled Graph Neural Networks,''
\textit{PyG Documentation v2.6.1}, 2024.
\url{https://pytorch-geometric.readthedocs.io/en/2.6.1/advanced/compile.html}

\bibitem{Lin2024}
N.~Lin, S.~Orfanoudakis, N.~O.~Cardenas, J.~S.~Giraldo, and P.~P.~Vergara,
``PowerFlowNet: Power flow approximation using message passing Graph Neural
Networks,''
\textit{Int.\ J.\ Electr.\ Power Energy Syst.}, vol.~160, p.~110112, 2024.
DOI:~10.1016/j.ijepes.2024.110112.
arXiv:~2311.03415.

\bibitem{Wang2021GPU}
Z.~Wang, S.~Wende-von Berg, and M.~Braun,
``Fast parallel Newton--Raphson power flow solver for large number of system
calculations with CPU and GPU,''
\textit{Sustain.\ Energy Grids Netw.}, vol.~27, p.~100488, 2021.
DOI:~10.1016/j.segan.2021.100488.
arXiv:~2101.02270.

\bibitem{Conrad2026}
T.~Conrad, C.~Kim, J.~J\"{a}ger, A.~Maier, and S.~Bayer,
``Impact of training dataset size for ML load flow surrogates,''
arXiv preprint arXiv:2602.19667, Feb.\ 2026.

\bibitem{Ghamizi2024}
S.~Ghamizi, J.~Cao, A.~Ma, and P.~Rodriguez,
``PowerFlowMultiNet: Multigraph neural networks for unbalanced three-phase
distribution systems,''
\textit{IEEE Trans.\ Power Syst.}, 2024.
DOI:~10.1109/TPWRS.2024.3465088.
arXiv:~2403.00892.

\bibitem{Shamseldein2025a}
M.~Shamseldein,
``A hybrid GNN-LSE method for fast, robust, and physically-consistent AC
power flow,''
arXiv preprint arXiv:2510.22020, Oct.\ 2025.

\bibitem{Kim2026}
C.~Kim, T.~Conrad, R.~Karim, J.~Oelhaf, D.~Riebesel, T.~Arias-Vergara,
A.~Maier, J.~J\"{a}ger, and S.~Bayer,
``Physics-informed GNN for medium-high voltage AC power flow with edge-aware
attention and line search correction operator,''
\textit{ICASSP 2026} (accepted).
arXiv:~2509.22458.

\bibitem{Bottcher2023}
L.~B\"{o}ttcher et al.,
``Solving AC power flow with graph neural networks under realistic
constraints,''
in \textit{Proc.\ 2023 IEEE Belgrade PowerTech}, pp.~1--7, 2023.
DOI:~10.1109/PowerTech55446.2023.10202246.
arXiv:~2204.07000.

\bibitem{Shamseldein2025b}
M.~Shamseldein,
``A hybrid GNN-IZR framework for fast and empirically robust AC power flow
analysis in radial distribution systems,''
arXiv preprint arXiv:2510.04264, Oct.\ 2025.

\bibitem{Donon2020}
B.~Donon, R.~Clément, B.~Donnot, A.~Marot, I.~Guyon, and M.~Schoenauer,
``Neural networks for power flow: Graph neural solver,''
\textit{Electr.\ Power Syst.\ Res.}, vol.~189, p.~106547, 2020.
DOI:~10.1016/j.epsr.2020.106547.

\end{thebibliography}

\end{document}
